% concatenated from timelike-conjugate-points.tex
\documentclass[a4paper,11pt]{article}
\usepackage[utf8x]{inputenc}

\usepackage[title]{appendix}

\usepackage{a4wide}
\usepackage{amssymb}
\usepackage{amsmath}
\usepackage{amsthm}
\usepackage{mathtools}
\usepackage[mathscr]{euscript}
\usepackage[normalem]{ulem}
\usepackage[colorlinks=true]{hyperref}
\usepackage[usenames,dvipsnames]{xcolor}
\usepackage{appendix}

\usepackage{graphicx}

\numberwithin{equation}{section}

\theoremstyle{plain}

\newtheorem{Definition}{Definition}[section]
\newtheorem{Construction}[Definition]{Construction}
\newtheorem{Theorem}[Definition]{Theorem}
\newtheorem{Proposition}[Definition]{Proposition}
\newtheorem{Lemma}[Definition]{Lemma}
\newtheorem{Corollary}[Definition]{Corollary}


\theoremstyle{definition}
\newtheorem{Remark}[Definition]{Remark}
\newtheorem{Example}[Definition]{Example}

% \newtheorem{Remark}[Definition]{Remark}
% \newtheorem{Example}[Definition]{Example}

\newcommand\red[1]{{\color{red} {#1}}}

\newcommand{\R}{\mathbb R}
\newcommand{\N}{\mathbb N}
\newcommand{\Fl}{\text{Fl}}
\newcommand{\Vol}{\mathrm{Vol}}
\newcommand{\Inj}{\mathrm{Inj}}
\newcommand{\grad}{\mathrm{grad}}
\newcommand{\pr}{\mathrm{pr}}
\newcommand{\fTCut}{\mathrm{1stCut}_T^+}
\newcommand{\fCCut}{\mathrm{1stCut}_C^+}
\newcommand{\fNCut}{\mathrm{1stCut}_N^+}

\DeclareMathOperator{\Geo}{Geo}
\DeclareMathOperator{\TGeo}{TGeo}

%mike:
\providecommand{\todoo}[1]{\P \textbf{#1} \P}
\newcommand{\eps}{\varepsilon}
%\newcommand{\vphi}{\varphi}
\usepackage[xcolor]{changebar}
\cbcolor{blue}
\newcommand{\Ric}{\mathrm{Ric}}
\newcommand{\sff}{\mathrm{II}} %second fundamental form
\newcommand{\gec}{{\check g_\eps}}
\newcommand{\geck}{{\check g_{\eps_j}}}
%%%\newcommand{gameck}{{\gamma_{\eps_j}}} I removed the backslash to make overleaf nicer
\newcommand{\comp}{\Subset}
\newcommand{\X}{\mathfrak{X}}
\newcommand{\coo}{\mathcal{C}^{1,1}}
\newcommand{\ctw}{\mathcal{C}^{2}}
\newcommand{\sse}{\subseteq}
\newcommand{\conv}{\mathbf{k}}
\newcommand{\edge}{\mathrm{edge}}
\newcommand{\cone}{\mathcal{C}}
\newcommand{\linfloc}{L^\infty_{\mathrm{loc}}}
\newcommand{\D}{\mathcal{D}}
\newcommand{\diag}{\mathrm{diag}}
\newcommand{\lamin}{\lambda_{\text{min}}}
\newcommand{\Dpk}{\mathcal{D}'{}^{(k)}}
\newcommand{\Dpo}{\mathcal{D}'{}^{(1)}}
\newcommand{\trsm}{\mathcal{T}^r_s(M)}
\newcommand{\ltl}{L^2_{\mathrm{loc}}}
\newcommand{\lpl}{L^p_{\mathrm{loc}}}
\newcommand{\Cloc}{C^0_{\mathrm{loc}}}
\newcommand{\Celoc}{C^1_{\mathrm{loc}}}
\newcommand{\Om}{\Omega}
%\newcommand{\miketodo}[1]{todo[size=\tiny,color=cyan]{#1}} %other colors available: e.g.:red, green, blue, yellow, lightgrey, teal, lime, ...
%\newcommand{\rotodo}[1]{todo[size=\tiny,color=teal]{#1}} %
%\newcommand{\btodo}[1]{todo[size=\tiny,color=yellow]{#1}}
\newcommand{\tcb}[1]{\textcolor{blue}{#1}}
\newcommand{\tcr}[1]{\textcolor{green}{#1}}
\newcommand{\pt}{\partial}

\newcommand{\F}{\mathcal F}
\newcommand{\G}{\mathcal G}
\newcommand{\HH}{\mathcal H}

\newcommand{\diamfin}{\mathrm{diam_{fin}}}
\newcommand{\lm}[1]{\mathbb{L}^2(#1)}


%James
\newcommand{\cT}{{\mathcal{T}}}

%Melanie
\usepackage[inline]{enumitem}
\newcommand{\enumlabelformat}{\roman}
\newcommand{\enumlabelfont}[1]{#1}
\newlength{\thelabelsep}
\setlength{\thelabelsep}{4pt}

\setlist{labelsep=\thelabelsep}
\setlist[enumerate,1]{font=\enumlabelfont,label=(\enumlabelformat*),leftmargin=2.5em}
\setlist[itemize]{leftmargin=2.5em,label=$-$}

\newcounter{inlineenum}
\renewcommand{\theinlineenum}{\enumlabelformat{inlineenum}}
% use this to enumerate in remarks etc.
\newenvironment{inlineenum}
 {\setcounter{inlineenum}{0}%
  \renewcommand{\item}{\refstepcounter{inlineenum}{(\theinlineenum)\hspace{\thelabelsep}}}
 }
 {\ignorespacesafterend}

\newcommand\jp{\partial J^{+}(S)}
\newcommand\ep{E^{+}(S)}
\newcommand\dpep{D^{+}\left(E^{+}(S)\right)}
\newcommand\dpjp{D^{+}\left(\partial J^{+}(S)\right)}
\newcommand\hpjp{H^{+}\left(\partial J^{+}(S)\right)}
\newcommand\hpep{H^{+}\left(E^{+}(S)\right)}
\newcommand{\blue[1]}{{#1}}
%\newcommand{\R}{\mathbb{R}}
%providecommand{todo}[1]{\P \textbf{#1} \P}

\let\epsilon\varepsilon
\let\phi\varphi

%argam
\newcommand{\diam}{\mathrm{diam}}

%tobias
\newcommand{\fel}[1]{{\color{blue}{#1}}}
\newcommand{\tb}[1]{{\color{purple}{#1}}}
\newcommand{\argam}[1]{{\color{green!70!black}{#1}}}
\newcommand{\did}[1]{{\color{yellow!70!black}{#1}}}

%clemens
\newcommand{\LpLS}{Lorentzian pre-length space }
\newcommand{\LpLSn}{Lorentzian pre-length space}
\newcommand{\ma}{\measuredangle}

%felix
%widebar command (i prefer to use this for most things over bar)
\makeatletter
\let\save@mathaccent\mathaccent
\newcommand*\if@single[3]{%
  \setbox0\hbox{${\mathaccent"0362{#1}}^H$}%
  \setbox2\hbox{${\mathaccent"0362{\kern0pt#1}}^H$}%
  \ifdim\ht0=\ht2 #3\else #2\fi
  }
%The bar will be moved to the right by a half of \macc@kerna, which is computed by amsmath:
\newcommand*\rel@kern[1]{\kern#1\dimexpr\macc@kerna}
%If there's a superscript following the bar, then no negative kern may follow the bar;
%an additional {} makes sure that the superscript is high enough in this case:
\newcommand*\widebar[1]{\@ifnextchar^{{\wide@bar{#1}{0}}}{\wide@bar{#1}{1}}}
%Use a separate algorithm for single symbols:
\newcommand*\wide@bar[2]{\if@single{#1}{\wide@bar@{#1}{#2}{1}}{\wide@bar@{#1}{#2}{2}}}
\newcommand*\wide@bar@[3]{%
  \begingroup
  \def\mathaccent##1##2{%
%Enable nesting of accents:
    \let\mathaccent\save@mathaccent
%If there's more than a single symbol, use the first character instead (see below):
    \if#32 \let\macc@nucleus\first@char \fi
%Determine the italic correction:
    \setbox\z@\hbox{$\macc@style{\macc@nucleus}_{}$}%
    \setbox\tw@\hbox{$\macc@style{\macc@nucleus}{}_{}$}%
    \dimen@\wd\tw@
    \advance\dimen@-\wd\z@
%Now \dimen@ is the italic correction of the symbol.
    \divide\dimen@ 3
    \@tempdima\wd\tw@
    \advance\@tempdima-\scriptspace
%Now \@tempdima is the width of the symbol.
    \divide\@tempdima 10
    \advance\dimen@-\@tempdima
%Now \dimen@ = (italic correction / 3) - (Breite / 10)
    \ifdim\dimen@>\z@ \dimen@0pt\fi
%The bar will be shortened in the case \dimen@<0 !
    \rel@kern{0.6}\kern-\dimen@
    \if#31
      \overline{\rel@kern{-0.6}\kern\dimen@\macc@nucleus\rel@kern{0.4}\kern\dimen@}%
      \advance\dimen@0.4\dimexpr\macc@kerna
%Place the combined final kern (-\dimen@) if it is >0 or if a superscript follows:
      \let\final@kern#2%
      \ifdim\dimen@<\z@ \let\final@kern1\fi
      \if\final@kern1 \kern-\dimen@\fi
    \else
      \overline{\rel@kern{-0.6}\kern\dimen@#1}%
    \fi
  }%
  \macc@depth\@ne
  \let\math@bgroup\@empty \let\math@egroup\macc@set@skewchar
  \mathsurround\z@ \frozen@everymath{\mathgroup\macc@group\relax}%
  \macc@set@skewchar\relax
  \let\mathaccentV\macc@nested@a
%The following initialises \macc@kerna and calls \mathaccent:
  \if#31
    \macc@nested@a\relax111{#1}%
  \else
%If the argument consists of more than one symbol, and if the first token is
%a letter, use that letter for the computations:
    \def\gobble@till@marker##1\endmarker{}%
    \futurelet\first@char\gobble@till@marker#1\endmarker
    \ifcat\noexpand\first@char A\else
      \def\first@char{}%
    \fi
    \macc@nested@a\relax111{\first@char}%
  \fi
  \endgroup
}
\makeatother
%comparison points
\newcommand*{\bx}{\bar{x}}
\newcommand*{\by}{\bar{y}}
\newcommand*{\bz}{\bar{z}}
\newcommand*{\bv}{\bar{v}}
\newcommand*{\bw}{\bar{w}}
\newcommand*{\bp}{\bar{p}}
\newcommand*{\bq}{\bar{q}}
\newcommand*{\br}{\bar{r}}
\newcommand*{\ba}{\bar{a}}
\newcommand*{\bb}{\bar{b}}

%other comparison stuff
\newcommand*{\btau}{\bar{\tau}}

\newcommand{\mike}[1]{\textcolor{blue}{#1}}
\newcommand{\jamesie}[1]{\textcolor{magenta}{#1}}
\newcommand\LpLSmaths{\ensuremath{(X, d, \ll, \leq, \tau)}}

%\renewcommand{\todo}[1]{}
%\renewcommand{\todoo}[1]{}




\newcommand{\g}{\gamma}
\newcommand{\s}{\sigma}
\newcommand{\vp}{\varphi}
\newcommand{\ps}{\psi}
\newcommand{\hvp}{\hat{\varphi}}
\newcommand{\hps}{\hat{\psi}}



%----- End felix commands -----

\title{Timelike conjugate points in Lorentzian length spaces}

\author{James D.E.\ Grant\footnote{School of Mathematics and Physics, University of Surrey, UK, j.grant@surrey.ac.uk}\\
Michael Kunzinger\footnote{Department of Mathematics, University of Vienna, Oskar-Morgenstern-Platz 1, 1090 Wien, Austria, \newline michael.kunzinger@univie.ac.at, b.schinnerl@gmail.com, roland.steinbauer@univie.ac.at},\\
Argam Ohanyan\footnote{Department of Mathematics, University of Toronto, 45 St.\ George Street, M5S 2E5 Toronto, Ontario, Canada, argam.ohanyan@utoronto.ca},\\
Yasmin Schinnerl${}^\dagger$,\\
Roland Steinbauer${}^\dagger$
}

\begin{document}

\date{\today}

%\date{Received: date /Accepted: date}

\maketitle

\begin{abstract}

We study notions of conjugate points along timelike geodesics in the synthetic setting of Lorentzian (pre-)length spaces, inspired by earlier work for metric spaces by Shankar--Sormani. After preliminary considerations on convergence of timelike and causal geodesics, we introduce and compare one-sided, symmetric, unreachable and ultimate conjugate points along timelike geodesics. We show that all such notions are compatible with the usual one in the smooth (strongly causal) spacetime setting. As applications, we prove a timelike Rauch comparison theorem, as well as a result closely related to the recently established Lorentzian Cartan--Hadamard theorem by Er\"{o}s--Gieger. In the appendix, we give a detailed treatment of the Fr\'{e}chet distance on the space of non-stopping curves up to reparametrization, a technical tool used throughout the paper.

\vspace{1em}

\noindent
\emph{Keywords:} Timelike conjugate point, geodesic, timelike Rauch comparison
\medskip

\noindent
\emph{MSC2020:} 53B30, 53C22, 53C23

\end{abstract}

\tableofcontents


\section{Introduction}\label{section: introduction}

Lorentzian (pre-)length spaces were introduced by Kunzinger--S\"{a}mann~\cite{kunzinger2018lorentzian} as a way to study Lorentzian and spacetime geometry using only the most essential structures, akin to metric geometry in positive definite signature. Since its introduction, the synthetic theory of spacetimes has experienced rapid development, and is a very active field of research today.

Geodesics (i.e., local distance realizers) in Lorentzian pre-length spaces were introduced and studied by Grant--Kunzinger--S\"{a}mann~\cite{grant2019inextendibility}. This notion turned out to be relevant in the study of hyperbolic angles by Beran--S\"{a}mann~\cite{beran2022angles} and the recently established synthetic version of the sectional curvature variant of Bartnik's splitting conjecture by Flores--Herrera--Solis~\cite{flores2024low}.

In this article, we are interested in studying notions of conjugate points along timelike geodesics in Lorentzian (pre-)length spaces, 
following the fundamental contribution of Shankar--Sormani~\cite{shankar2009conjugate} on conjugate points in metric spaces. (See also the earlier work of Rinow~\cite{rinow1961innere}.) In the smooth Riemannian setting, Warner's detailed study~\cite{warner1965conjugate} of the conjugate locus sheds light on the exact behavior of the exponential map and the type of singularities that can occur. The results of Warner were extended to semi-Riemannian manifolds by Szeghy~\cite{szeghy2008conjugate}.

In Lorentzian and spacetime geometry, conjugate points are fundamental in the study of various comparison-geometric results, such as the singularity theorems of Hawking and Penrose. Our aim in this paper is to initiate the study of conjugate points in the synthetic setting for Lorentzian geometry, i.e. Lorentzian pre-length spaces. 

Our article is structured as follows: In Section~\ref{Section: Geodesics}, we first recall some basics on Lorentzian pre-length spaces used throughout (Subsection~\ref{Subsection: preliminaries on LpLS}). We then study geodesics and their convergence properties in Subsection~\ref{subsection: convergenceofgeodesics}. Due to the lack of arclength parametrization in the generality that we work with, we need to consider equivalence classes of geodesics up to suitable reparametrizations. To this end, in Appendix~\ref{Section: appendix frechet}, we give a treatment of the~\emph{Fr\'{e}chet distance\/} on such equivalence classes of non-stopping curves on a metric spaces that is suitable for our purposes. Our main result in this context is that the set of (equivalence classes of) causal geodesics $\Geo(X)$ is closed in the set of non-stopping curves with respect to the Fr\'{e}chet distance (see Proposition~\ref{Proposition: TopologyofGeoX}). 

In Section~\ref{Section: Timelikeconjugatepoints}, following the terminology of Shankar--Sormani~\cite{shankar2009conjugate}, we introduce and study various notions of conjugate points along timelike geodesics: One-sided and symmetric conjugate points in Subsection~\ref{Subsection: symmetric conjugate points}, families of timelike geodesics in Subsection~\ref{subsection: Families of timelike geodesics}, unreachable and ultimate conjugate points in Subsection~\ref{Subsection: unreachable and ultimate}, as well as some notions of cut points in Subsection~\ref{Subsection: causalcutloci}. We rely on Szeghy's~\cite{szeghy2008conjugate} generalizations of Warner's~\cite{warner1965conjugate} results to show compatibility of these notions with the smooth (strongly causal) spacetime setting.

Section~\ref{Section: applications} deals with two applications of these notions: a timelike Rauch comparison theorem (Theorem~\ref{Theorem: Timelike Rauch}) and a result similar to the Cartan--Hadamard Theorem (Theorem~\ref{Theorem: cartan hadamard adjacent}), which establishes the nonexistence of a certain kind of conjugate points along sufficiently short geodesics in spaces with upper timelike curvature bounds. The main tools that we use (aside from those that we develop here) is the stacking principle of Beran--Ohanyan--Rott--Solis~\cite{beran2023splitting} and the synthetic Lorentzian Cartan--Hadamard theorem of Er\"{o}s--Gieger~\cite{Eroes-Gieger25}.

Finally, in Section~\ref{section: Outlook}, we give an outlook on related avenues of future research.

\section{Geodesics in Lorentzian pre-length spaces}
\label{Section: Geodesics}

\subsection{Preliminaries on Lorentzian pre-length spaces}
\label{Subsection: preliminaries on LpLS}

In this section, we collect a few basics on Lorentzian (pre-)length spaces that will be required throughout. We refer to Kunzinger--S\"{a}mann~\cite{kunzinger2018lorentzian} for full details.

\begin{Definition}[Lorentzian (pre-)length spaces]
A~\emph{Lorentzian pre-length space} is a quintuple $(X,d,\ll,\leq,\tau)$, where
\begin{enumerate}
    \item $(X,d)$ is a metric space; 
    \item $\ll, \leq$ are transitive relations (the~\emph{chronological\/} resp.\ \emph{causal relations}), $\leq$ is in addition reflexive and contains $\ll$; 
    \item $\tau \colon X^2 \to [0,\infty]$ is a lower semicontinuous function that vanishes outside of the causal relation $\leq$, the timelike relation $\ll$ coincides with $\{\tau > 0\}$ and for all $x \leq y \leq z$,
    \begin{equation}
        \tau(x,z) \geq \tau(x,y) + \tau(y,z).
    \end{equation}
\end{enumerate}
\end{Definition}

We employ standard notation $I^\pm(A), J^\pm(A)$ for chronological resp.\ causal futures/pasts of a subset $A \subseteq X$. We write $I(A,B) \coloneqq I^+(A) \cap I^-(B)$, similarly $J(A,B)$, for $A, B \subseteq X$. A subset $C \subseteq X$ is said to be~\emph{causally convex\/} if $J(x,y) \subseteq C$ whenever $x,y \in C$. A Lorentzian pre-length space is~\emph{chronological\/} if $\ll$ is irreflexive and~\emph{causal\/} if $\leq$ is a partial order.

\begin{Definition}[Timelike/causal curves and $\tau$-length]
\label{def:TLcausalcurves}
A~\emph{future-directed timelike (resp.\ causal) curve\/} (past-directed curves are defined analogously) in a Lorentzian pre-length space $X$ is a nowhere constant map $\gamma \colon I=[a,b]\to X$ (i.e., $\gamma$ is not constant on any non-trivial subinterval of $I$) that
is Lipschitz (w.r.t.\ $d$) and such that for all $t_1, t_2 \in I$ with $t_1<t_2$ we have $\gamma(t_1) \ll \gamma(t_2)$ (resp.) $\gamma(t_1)\le \gamma(t_2)$ (causal curves on unbounded intervals can be defined analogously). Its~\emph{$\tau$-length\/} is defined as 
\[
L_\tau(\gamma) = \inf \left\{ \sum_{i=0}^{N-1} \tau(\gamma(t_i),\gamma(t_{i+1})) \colon 
N\in \N,\ a=t_0 < t_1 < \dots < t_N = b \right\}.
\]
A reparametrization of a future-directed causal curve $\gamma \colon [a,b] \to X$ is a future-directed causal curve $\lambda \colon [c,d] \to X$ such that $\gamma=\lambda\circ \varphi$, where $\varphi \colon [a,b]\to [c,d]$ is surjective and strictly monotonically increasing (hence also continuous). A future-directed causal curve $\gamma \colon [a,b] \to X$ is called a~\emph{maximizer\/} if $L_\tau(\gamma) = \tau(\gamma(a),\gamma(b))$.
\end{Definition}

In particular, the reparametrized curve $\lambda\circ\varphi$ above is required to itself be Lipschitz. The $\tau$-length of a causal curve is invariant under reparametrizations (\cite[Lem.\ 2.18]{kunzinger2018lorentzian}). We note that causal curves were defined in~\cite{kunzinger2018lorentzian} as non-constant maps rather than nowhere constant. We choose the current convention for two main reasons, one being that the restriction of a causal curve to any non-trivial subinterval of its domain should again be causal, and the other that causal curves in the standard smooth spacetime setting are indeed nowhere constant. The Fr\'{e}chet distance as detailed in Appendix~\ref{Section: appendix frechet} then applies to causal curves.

We also note that all causal curves are assumed to be future-directed unless explicitly stated otherwise.

\begin{Definition}[Model spaces]
\label{Def: model spaces}
For integers $0 \le m \le n$, denote by $\mathbb{R}_{m}^{n}$ the vector space $\mathbb{R}^{n}$ with the inner product $b(v,w) \coloneqq -\sum_{i=1}^{m}v_{i}w_{i} + \sum_{i=m+1}^{n}v_{i}w_{i},$ where $v=(v_1, \dots, v_n)$ and $w=(w_1, \dots ,w_n)$.\par
Let $K \in \mathbb{R}$. The~\emph{Lorentzian $K$-planes\/} or the~\emph{comparison spaces of constant curvature $K$\/} are the following:
\begin{itemize}
    \item For $K>0$, $\mathbb{L}^{2}(K)$ is the universal cover of $\left\{ v \in \mathbb{R}_{1}^{3} : b(v,v) = \frac{1}{K^{2}} \right\}$.
    \item For $K<0$, $\mathbb{L}^{2}(K)$ is the universal cover of $\left\{v \in \mathbb{R}_{2}^{3} : b(v,v) = -\frac{1}{K^{2}} \right\}$.
    \item For $K=0$, $\mathbb{L}^{2}(0) \coloneqq \mathbb{R}_{1}^{2}$ is the two-dimensional Minkowski spacetime.
\end{itemize}
The timelike diameter of $\lm{K}$ is $D_K = \frac{\pi}{\sqrt{-K}}$ for $K < 0$ and
$D_K = \infty$ for $K \ge 0$. We will denote the time separation of $\mathbb{L}^2(K)$ by $\bar{\tau}$.
\end{Definition}

\begin{Definition}[Alexandrov topology, strong causality]
The~\emph{Alexandrov topology\/} on a Lorentzian pre-length space $X$ is defined as the topology generated by the subbase
\begin{equation}
    \{I(x,y) : x,y \in X\}.
\end{equation}
We say $X$ is~\emph{strongly causal\/} if the Alexandrov topology agrees with the metric topology on $X$.
\end{Definition}

Following~\cite{BKR24}, we call a Lorentzian pre-length space~\emph{regular\/} if every maximizer between two timelike related points is timelike, i.e., contains no null subsegment. Moreover, a Lorentzian pre-length space is called~\emph{causally path-connected\/} if any two distinct causally related points are joined by a causal curve and any two timelike related points by a timelike curve.

\begin{Remark}[Localizability]
A~\emph{localizable\/} Lorentzian pre-length space is a Lorentzian pre-length space equipped with well-behaved local time separation functions, see~\cite[Def.\ 3.16]{kunzinger2018lorentzian} for the precise definition. It is called~\emph{regularly localizable\/} if the maximizers w.r.t.\ the local time separations (which exist by definition of localizability) are required to be  
timelike if they connect timelike related points. Regularly localizable Lorentzian pre-length spaces are regular, cf.\ \cite[Thm.\ 3.18]{kunzinger2018lorentzian}.
\end{Remark}

\begin{Proposition}[Upper semicontinuity of $L_\tau$, Prop.\ 3.17 in~\cite{kunzinger2018lorentzian}]
Let $X$ be a localizable and strongly causal Lorentzian pre-length space. Then $L_\tau$ is upper semicontinuous in the following sense: If $\gamma_n \colon [a,b] \to X$ is a sequence of causal curves converging $d$-uniformly to a causal curve $\gamma \colon [a,b] \to X$, then
\begin{equation}
    L_\tau(\gamma) \geq \limsup_{n \to \infty} L_\tau(\gamma_n).
\end{equation}
\end{Proposition}

For our definition of timelike (sectional) curvature bounds, we follow the conventions of Beran--Kunzinger--Rott~\cite[Def.\ 3.1]{BKR24}.

\begin{Definition}[Timelike curvature bounds]
We say a Lorentzian pre-length space $X$ has~\emph{timelike curvature bounded from above (resp.\ below)\/} by $K \in \R$ if every point has a neighborhood $U$ such that
\begin{enumerate}
    \item the set $U^2 \cap \tau^{-1}([0,D_K))$ is open in $X^2$ and the restriction of $\tau$ to this set is continuous and finite-valued;
    \item whenever $x \leq y$ with $\tau(x,y) < D_K$ and $x,y \in U$, there exists a maximizer in $U$ from $x$ to $y$;
    \item whenever $\Delta(x,y,z)$ is a timelike triangle in $U$ (i.e., $x \ll y \ll z$ with timelike maximizers running between them) and $\Delta(\bar{x},\bar{y},\bar{z})$ is the comparison triangle with equal side lengths in $\mathbb{L}^2(K)$, for any $p,q \in \Delta(x,y,z)$ and corresponding points $\bar{p},\bar{q} \in \Delta(\bar{x},\bar{y},\bar{z})$ (i.e., points with equal time separation to the endpoints of their sides as $p,q$), we have
    \begin{equation}
        \tau(p,q) \geq \bar{\tau}(\bar{p},\bar{q}) \quad (\text{resp.\ }\leq).
    \end{equation}
\end{enumerate}
\end{Definition}

We will mainly rely on consequences of timelike curvature bounds, specifically the stacking principle~\cite[Prop.\ 2.42]{beran2023splitting}, which describes how the stacking of a subdivision of comparison triangles stack in the model space, as well as the Lorentzian Cartan--Hadamard theorem and related results due to Er\"{o}s--Gieger~\cite{Eroes-Gieger25}. We refer to these works for details.

\subsection{Convergence of geodesics}\label{subsection: convergenceofgeodesics}

In this section, we establish results on the convergence of geodesics in Lorentzian (pre-)length spaces. We begin by recalling the notion of geodesic introduced in~\cite[Def.\ 1.7]{beran2022angles}: 
\begin{Definition}\label{def:geodesics}
A future-directed causal curve $\gamma \colon I \to X$ is a~\emph{geodesic\/} if for each $t \in I$ there
exists a neighborhood $[a,b]$ of $t$ (i.e., $a < t < b$, allowing for equality at the
endpoints of $I$) such that $\gamma|_{[a,b]}$ is a maximizer, i.e., $\tau(\gamma(a), \gamma(b)) =
L_\tau (\gamma|_{[a,b]})$.
\end{Definition}

Since the $\tau$-length of a causal curve $\gamma$ is parametrization-invariant, any reparametrization of a geodesic $\gamma$ is also a geodesic, and we shall henceforth consider geodesics as equivalence classes of reparametrizations. We emphasize that whenever we speak of "uniqueness" in the context of geodesics, we always mean "up to reparametrization".

\begin{Definition}[Finite $\tau$-measurability] A Lorentzian pre-length space $(X,d,\ll,\leq,\tau)$ is~\emph{finitely $\tau$-measurable\/} if every causal curve segment $\gamma \colon [a,b] \to X$ has finite $\tau$-length.
\end{Definition}

\begin{Lemma}
Let {\LpLSmaths} be a Lorentzian pre-length space that is localizable or has a timelike curvature bound (from above or below). Then it is finitely $\tau$-measurable.
\end{Lemma}
\begin{proof}
The first claim follows from the definition of localizability, cf.\ \cite[Def.\ 3.16]{kunzinger2018lorentzian} for the latter. On the other hand, if $X$ has a timelike curvature bound, then $\tau$ is continuous and finite-valued on the product with itself of a comparison neighborhood of any point in $X$. Covering a given causal curve segment with finitely many such neighborhoods establishes finite $\tau$-measurability.
\end{proof}

\begin{Remark} \label{Remark: geodesicsvssolutionsofgeodeq} 
In~\cite[Def.\ 4.1]{grant2019inextendibility}, geodesics were defined as future-directed timelike (resp.\ causal, null) curves $\gamma \colon I \to X$ on localizing Lorentzian pre-length spaces as follows: for every $t_0 \in I$ there exists a localizing neighborhood $\Omega$ of $\gamma(t_0)$ and a relatively closed interval $J \subseteq I$ with $t_0$ in its relative interior, such that $\gamma(J) \subseteq \Omega$ and $\gamma|_J$ is maximizing in $\Omega$ with respect to the local time separation function $\omega$. Thus, such curves are geodesics in the sense of Definition~\ref{def:geodesics} within each `local' Lorentzian pre-length space $\Omega$ as above. Moreover, in strongly causal localizable Lorentzian pre-length spaces both definitions agree by the proof of~\cite[Lem.\ 4.3]{grant2019inextendibility}.
\end{Remark}

\begin{Remark}[On geodesics in low causality]
The definition of causal geodesic we gave as a local $\tau$-maximizer becomes meaningless in very low causality. For example, consider the Lorentz cylinder $X = S^1 \times \R^n$ with the metric $-d\theta^2 + \sum_{i=1}^n (dx^i)^2$, then $\tau = +\infty$ on the timelike relation $\ll$. In particular, the $\tau$-length of every timelike curve segment is $+\infty$, hence every such curve is maximizing and thus (in particular) a geodesic.
\end{Remark}


An important class of examples of Lorentzian (pre-) length spaces is furnished by smooth manifolds $M$ equipped with a continuous and causally plain metric $g$ (cf.\ \cite[Prop.\ 5.8]{kunzinger2018lorentzian}). To compare the notion of geodesics in Lorentzian manifolds with that of Definition~\ref{def:geodesics}, let us now extend the validity of~\cite[Prop.\ 2.32]{kunzinger2018lorentzian}
from smooth Lorentzian metrics to a large class of continuous ones. 

\begin{Proposition}[$L_g$ vs.\ $L_\tau$]
\label{Proposition: Lg vs Ltau}
Let $(M,g)$ be a spacetime such that $g$ is continuous, causally plain and strongly causal. Then the $g$- and $\tau$-lengths of every causal curve $\gamma \colon [a, b] \to (M,g)$ agree:
\[
L_g(\gamma) = \int_a^b \sqrt{-g(\dot\gamma,\dot\gamma)} \,dt = L_\tau(\gamma).
\]
\end{Proposition}
\begin{proof}
To prove this we employ a sequence of smooth Lorentzian metrics $g_k$ with wider lightcones than $g$, approximating $g$ as constructed in~\cite[Lem.\ A.1]{McCann_Saemann22}. Then for the corresponding time separation functions $\tau_k$ of $g_k$ we have that $\tau$ is the locally uniform and non-increasing limit of the $\tau_k$ (see~\cite[Lem.\ A.2]{McCann_Saemann22}). In particular, $L_{\tau_k}(\gamma) \ge L_\tau(\gamma)$ for all $k$. Let $\eps>0$ and choose a partition $P=\{a=t_0<\dots<t_{N}=b\}$ of $[a,b]$ such that
\[
L_\tau(\gamma) \ge \sum_{i=1}^N \tau(\gamma(t_{i-1}),\gamma(t_i)) - \eps.
\]
By uniform convergence, there exists some $k_0=k_0(\eps,P)$ such that, for all $k\ge k_0$ we have
\[
\sum_{i=1}^N \tau(\gamma(t_{i-1}),\gamma(t_i)) \ge \sum_{i=1}^N \tau_k(\gamma(t_{i-1}),\gamma(t_i)) -\eps.
\]
Altogether,
\[
L_{\tau_k}(\gamma) \ge L_\tau(\gamma) \ge \sum_{i=1}^N \tau_k(\gamma(t_{i-1}),\gamma(t_i)) - 2\eps 
\ge L_{\tau_k}(\gamma) -2\eps.
\]
Consequently, $L_\tau(\gamma) = \lim_{k\to\infty} L_{\tau_k}(\gamma) = \lim_{k\to \infty} L_{g_k}(\gamma) =L_g(\gamma)$.
\end{proof}

\begin{Remark}[Comparing notions of geodesics]
It follows from Proposition~\ref{Proposition: Lg vs Ltau} that (in strongly causal spaces) geodesics in the sense of Definition~\ref{def:geodesics} will coincide with geodesics
in the sense of solutions of the geodesic equation as long as the regularity of the metric is sufficient to allow for a satisfactory
solution theory of these ODEs and, additionally, as long as such solutions are indeed local maximizers of the $g$-length-functional.
This certainly is the case for $g\in C^{1,1}$ (cf.\ \cite{minguzzi2015convex} or~\cite{kunzinger2014regularisation}). However, 
the Hartmann--Wintner-type spacetime
$M \coloneqq \R \times (-1,1) \times \R^n$ with metric $g \coloneqq -dt^2 + dx^2 + (1 - |x|^{\lambda})g_n$, where $\lambda \in (1, 2)$ and $g_n$ is the Euclidean metric on $\R^n$, is 
an example of a $C^{1,\lambda-1}$-Lorentzian metric where solutions of the geodesic equations fail to even locally maximize the $g$-length, cf.\ \cite[Ex.\ 2.4]{samann2018geodesics}. In any case, whenever we compare synthetic notions to classical notions in spacetimes we will assume strong causality in order to have a compatible notion of local maximizers.
\end{Remark}

\begin{Lemma}[Local maximization and causal character] \ 
\label{Lemma: Geodesicshavecausalcharacter}
Let $(X,d,\ll,\leq,\tau)$ be a Lorentzian pre-length space. If $X$ is regular, then causal geodesics in $X$ have a causal character, i.e.\ they are either timelike or null.
\end{Lemma}
\begin{proof} 
Given a causal geodesic $\gamma \colon [a, b] \to X$, choose a subdivision $a=t_0<\dots<t_N=b$ such that each $\gamma|_{[t_i,t_{i+1}]}$ is maximal,
and therefore has a causal character due to regularity of $X$. Since, for each $i$, there is a neighborhood of $t_i$ in which $\gamma$ is maximal,
and therefore has a causal character, any two consecutive segments of $\gamma$ have the same causal character.
\end{proof}


Before discussing the limit curve theorem and its implications for geodesics, we adapt some material from~\cite{hau2020causal} to our setting.

\begin{Definition}[Local relations]
Let $(X,d,\ll,\leq,\tau)$ be a Lorentzian pre-length space and $U \subseteq X$ open. We define the local timelike and causal relations on $U$ as follows: For $p,q \in U$ we say $p \ll_U q$ ($p \leq_U q$) if there exists a future-directed timelike (causal) curve $\gamma \colon [a, b] \to X$ with $\gamma(a) = p$, $\gamma(b) = q$ and $\gamma \left( [a, b] \right) \subseteq U$.
\end{Definition}

\begin{Lemma}
Let {\LpLSmaths} be a causally path-connected {\LpLSn}. Then, for any causally convex $V \subseteq X$, we have $\le_V \, = \left. \le \right|_{V \times V}$.
\end{Lemma}

\begin{proof}
If $p, q \in V$ with $p \le_V q$, then there exists a future-directed, causal curve from $p$ to $q$, hence $p \le q$, so $(p, q) \in \left. \le \right|_{V \times V}$.

Conversely, let $p, q \in V$ with $p \le q$. Since {\LpLSmaths} is causally path-connected, there exists a causal curve $\gamma \colon [a, b] \to X$ with $\gamma(a) = p$, $\gamma(b) = q$. Since $V$ is causally convex, $\gamma \left( [a, b] \right) \subseteq V$. Therefore, $p \le_V q$.
\end{proof}

\begin{Definition}[Local weak causal closedness]
A set $U \subseteq X$ is~\emph{weakly causally closed\/} if, for all sequences $\left( p_n \right)_n$ and $\left( q_n \right)_n$ in $U$ such that $p_n \le_U q_n$, for all $n \in \N$ with $p_n \to p \in U$ and $q_n \to q \in U$, one has $p \le_U q$. A Lorentzian pre-length space is~\emph{locally weakly causally closed\/} if every point has a weakly causally closed neighborhood.
\end{Definition}

Note that, unlike~\cite[Definitions~2.18 \& 2.19]{hau2020causal}, we do not require that {\LpLSmaths} be causally path-connected in our definition of $\le_U$ and, hence, of weakly causally closed. However, we have the following slightly sharpened version of~\cite[Prop.~2.21]{hau2020causal}:

\begin{Lemma}
\label{Lemma:2.21}
Let $\LpLSmaths$ be a strongly causal Lorentzian pre-length space with the property that, for any causally convex $V \subseteq X$, we have $\le_V \, = \left. \le \right|_{V \times V}$. Then, the following are equivalent: 
\begin{enumerate}
\item\label{221i} $\LpLSmaths$ is locally causally closed;%
\footnote{Recall from~\cite[Definition~3.4]{kunzinger2018lorentzian} that if {\LpLSmaths} is a {\LpLSn}, then a subset $U \subseteq X$ is~\emph{causally closed\/} if, for all sequences $(p_n)_n$, $(q_n)_n$ in $U$ such that $p_n \le q_n$ for all $n$  with $p_n \to p \in \bar{U}$, $q_n \to q \in \bar{U}$, we have $p \le q$. A {\LpLSn} is~\emph{locally causally closed\/} if every point has a causally closed neighborhood.}
\item\label{221ii} $\LpLSmaths$ is locally weakly causally closed.
\end{enumerate}
\end{Lemma}

\begin{proof}
{\ }\\[-5mm]
\begin{itemize}[fullwidth]
(i) $\Rightarrow$ (ii): Let $x \in X$ and $U$ a causally closed neighborhood of $x$. By strong causality, there exists a causally convex neighborhood $V$ of $x$ with $x \in V \subseteq U$. Let $(p_n)_n$ and $(q_n)_n$ be sequences in $V$ with $p_n \le_V q_n$ for all $n$, such that $p_n \to p \in V$ and $q_n \to q \in V$ as $n \to \infty$. Since $V$ is causally convex, $\le_V \, = \left. \le \right|_{V \times V}$, so $p_n \le q_n$. Therefore, since $U$ is causally closed and $p, q \in V \subseteq U$, we have $p \le q$. Since $p, q \in V$ and $\le_V \, = \left. \le \right|_{V \times V}$, it follows that $p \le_V q$. Hence, $V$ is a weakly causally closed neighborhood of $x$. 

\item[(ii) $\Rightarrow$ (i)]: Let $x \in X$ and $U$ a weakly causally closed neighborhood of $x$. Since $(X, d)$ is a metric space, it satisfies the $T_3$ separation condition. In combination with strong causality, this implies that there exists a causally convex neighborhood $V$ of $x$ with the property that $x \in V \subseteq \overline{V} \subseteq U$. Let $(p_n)_n$ and $(q_n)_n$ be sequences in $V$ with $p_n \le q_n$ for all $n$, such that $p_n \to p$ and $q_n \to q$, where $p, q \in \overline{V} \subseteq U$. Since $V$ is causally convex, we have $p_n \le_V q_n$ and therefore, since $V \subseteq U$, $p_n \le_U q_n$ for all $n \in \N$. Since $U$ is weakly causally closed, it follows that $p \le_U q$ and, therefore, $p \le q$. Therefore, $V$ is a causally closed neighborhood of $x$.
\qedhere
\end{itemize}
\end{proof}

Finally, as pointed out in~\cite[pp.~8]{hau2020causal}, the limit curve theorem~\cite[Thm.\ 3.7]{kunzinger2018lorentzian} holds with the assumption of local weak causal closedness.%

\begin{Theorem}[Limit curve theorem for segments]
\label{Theorem: limitcurvetheoremforsegments}
Let $(X,d,\ll,\leq,\tau)$ be a 
locally weakly causally closed Lorentzian pre-length space. Let $\gamma_n \colon [a, b] \to X$ be a uniformly Lipschitz sequence of future-directed causal curves. Suppose that the $\gamma_n$ are either
\begin{itemize}
    \item[(a)] all contained in a common compact set, or 
    \item[(b)] the metric $d$ is proper, and there is $t_0 \in [a, b]$ such that a subsequence of $\gamma_n(t_0)$ converges. 
\end{itemize}
Then a subsequence of the $\gamma_n$ converges $d$-uniformly to a Lipschitz curve $\gamma \colon [a, b] \to X$. If $\gamma$ is nowhere constant, then $\gamma$ is a future-directed causal curve. Moreover:
\begin{enumerate}
\item\label{LCi} If $\gamma_n(a) = p$ and $\gamma_n(b) = q$ for each $n \in \N$, with $p \neq q$, then $\gamma$ is a future-directed causal curve from $p$ to $q$; 
\item\label{LCii} If $X$ is, in addition, strongly causal and localizable and the $\gamma_n$ are limit maximizing, i.e.\ there exists a sequence $\varepsilon_n \to 0$ such that
\[
L_{\tau}(\gamma_n) \geq \tau(\gamma_n(a),\gamma_n(b)) - \varepsilon_n,
\]
then $\gamma$ is maximizing. In this case, $L_{\tau}(\gamma) = \lim_{k \to \infty} L_{\tau}(\gamma_{n_k})$, where $\gamma_{n_k}$ is the subsequence converging to $\gamma$.
\end{enumerate}
\end{Theorem}
\begin{proof}
From the proof of~\cite[Thm.~3.7]{kunzinger2018lorentzian}, we deduce that there exists a subsequence $\left( \gamma_{n_k} \right)_k$ of curves that converge uniformly to a Lipschitz continuous limit curve $\gamma$. Therefore, without loss of generality, we may assume that the sequence of curves $\left( \gamma_n \right)_n$ converges uniformly to $\gamma$. Since $X$ is locally weakly causally closed, for every $t \in [a, b]$ there exists an open, weakly causally closed neighborhood $U_t$ of $\gamma(t)$. Let $a \le t_1 < t_2 \le b$ be such that $\gamma(t_1), \gamma(t_2) \in U_t$. Then, by uniform convergence, there exists an $n_0 \in \N$ such that, for all $n \ge n_0$, we have $\gamma_n \!\left( \left[ t_1, t_2 \right] \right) \subseteq U_t$. Since $\gamma_n$ are causal, this implies that $\gamma_n(t_1) \le_{U_t} \gamma_n(t_2)$ for all $n \ge n_0$. The fact that $U_t$ is weakly causally closed then implies that $\gamma(t_1) \le_{U_t} \gamma(t_2)$ and, therefore, $\gamma(t_1) \le \gamma(t_2)$. The rest of the proof of part~\ref{LCi} proceeds as in the proof of~\cite[Lemma~3.6]{kunzinger2018lorentzian}. 


The claim~\ref{LCii} about the maximality of the limit is a consequence of lower semicontinuity of $\tau$ and upper semicontinuity of $L_{\tau}$, 
cf.\ the proof of~\cite[Thm.\ 2.23]{beran2023splitting}.
\end{proof}

In the context of geodesics, we will often be interested in Lorentzian pre-length spaces where the local maximization property is uniform. This motivates the following definitions (one of which was already introduced in~\cite[Def.\ 2.20]{beran2023splitting}).

\begin{Definition}[Locally maximizing]\label{def:locally-maximising}
A Lorentzian pre-length space $(X,d,\ll,\leq,\tau)$ is called~\emph{locally timelike (causally, null) maximizing}\footnote{Such spaces were referred to as~\emph{locally quasiuniformly maximizing} in~\cite[Def.\ 2.20]{beran2023splitting}.} if for each point $p \in X$ there exists a neighborhood $U$ of $p$ such that each timelike (causal, null) geodesic segment entirely contained in $U$ is maximizing.
\end{Definition}


If $X$ is locally causally maximizing, we simply say that it is locally maximizing.

\begin{Example}[$C^{1,1}$-spacetimes are locally maximizing]
Let $(M,g)$ be a strongly causal $C^{1,1}$-spacetime, then it is locally maximizing: By Remark~\ref{Remark: geodesicsvssolutionsofgeodeq}, the two available notions of timelike (null, causal) geodesics agree. Moreover, by \cite[Thm.\ 6]{minguzzi2015convex} local maximizers coincide with classical solutions of the geodesic equation.
Now let $p \in M$ and $U$ a convex neighborhood containing $p$ (cf.\ \cite{kunzinger2014regularisation,minguzzi2015convex}). Then 
for any two causally related points $q_1 \leq q_2$, $q_1,q_2\in U$, there is a unique causal geodesic connecting them which is entirely contained in $U$. Hence any causal geodesic in $U$ is already maximizing.
\end{Example}

\begin{Example}[Low regularity spacetimes which are not locally maximizing] The Hartman-Wintner-type spacetimes from \cite[Ex.\ 2.4]{samann2018geodesics}, which are $C^{1,\alpha}$ for $0<\alpha<1$, are not locally maximizing.
\end{Example}

\begin{Example}
    The Lorentzian constant curvature $K$ spaces, which are (the universal coverings of) de-Sitter space, Minkowski space, and anti-de Sitter space of the corresponding radius (cf.\ Definition~\ref{Def: model spaces}), are locally maximizing: Indeed in the case of $K \geq 0$ this is trivial since causal geodesics are maximizing throughout, and in the case $K < 0$ one can consider convex neighborhoods of a given point $x$ which avoid the timelike cut locus 
    $\{\tau(x,.) = \frac{\pi}{\sqrt{-K}}\}$.
\end{Example}

\begin{Definition}[Spaces of geodesics]
Let $(X,d,\ll,\leq,\tau)$ be a finitely $\tau$-measurable Lorentzian pre-length space. We denote by $\Geo(X)$ the set of equivalence classes of future-directed causal geodesic segments $\gamma \colon [0, 1] \to X$, where two such segments are called equivalent if and only if there exists a  strictly increasing and surjective reparametrization $\varphi \colon [0,1] \to [0,1]$ between them.%
\footnote{See Definition~\ref{def:TLcausalcurves}.} %
We denote by $\TGeo(X) \subseteq \Geo(X)$ the subset of equivalence classes of future-directed timelike geodesic segments. For convenience, we write $\gamma \in \Geo(X)$ for both the equivalence class and some representative. 

We consider the following metrics on $\Geo(X)$: $d_F$ is the restriction of the Fr\'{e}chet distance (see Appendix~\ref{Section: appendix frechet}), and
\begin{align}
    d_{\Gamma}(\gamma_1,\gamma_2) \coloneqq d_F(\gamma_1,\gamma_2) + |L_{\tau}(\gamma_1) - L_{\tau}(\gamma_2)|.
\end{align}
\end{Definition}

In the following result, we denote by $\mathrm{Const}(X)$ the set of all constant paths $[0,1]\to X$.
\begin{Proposition}[Topology of $\Geo(X)$]
\label{Proposition: TopologyofGeoX}
Let $(X,d,\ll,\leq,\tau)$ be a localizable, locally weakly causally closed, strongly causal, and locally maximizing Lorentzian pre-length space. Then $\Geo(X)$ is a closed subset of the set of nowhere constant curves from $[0,1]$ to $X$ w.r.t.\ the Fr\'{e}chet distance $d_F$. Moreover, the metric topologies of $d_F$ and $d_\Gamma$ coincide on $\Geo(X)$.
\begin{proof}
    Let $(\gamma_n) \subseteq \Geo(X)$ converge to the equivalence class of a nowhere constant path $\gamma \colon [0,1] \to X$. By definition of $d_F$, there exists a sequence of representatives (again called $\gamma_n$) that converge uniformly to $\gamma$ (chosen to be some representative). By~\cite[Lem.\ 3.6]{kunzinger2018lorentzian} (cf.\ also~\cite{hau2020causal}), since $\gamma$ is assumed to be nowhere constant, it is a future-directed causal curve. Given $t \in [0,1]$, let $U$ be a neighborhood containing $\gamma(t)$ such that each causal geodesic entirely contained in $U$ is maximizing. W.l.o.g.\ suppose $t \in (0,1)$ and suppose $\delta >0$ so that $\gamma([t-\delta,t+\delta]) \subseteq U$. By uniform convergence, $\gamma_n([t-\delta,t+\delta])\subseteq U$ for all large enough $n$. Since $L_{\tau}$ is upper semicontinuous, one may then argue just like in the case of the limit curve theorem (cf.\ Theorem~\ref{Theorem: limitcurvetheoremforsegments}) to conclude that $\gamma([t-\delta,t+\delta])$ is maximizing, and in fact $\lim_n L_{\tau}(\gamma_n|_{[t-\delta,t+\delta]}) = L_{\tau}(\gamma|_{[t-\delta,t+\delta]})$ (cf.\ again the proof of~\cite[Thm.\ 2.23]{beran2023splitting}). Since $t$ was arbitrary, we conclude that $\gamma \in \Geo(X)$. Due to the convergence of Lorentzian lengths we just pointed out, we also see that $d_{F}$ and $d_{\Gamma}$ induce the same topology on $\Geo(X)$.
\end{proof}
\end{Proposition}



\medskip
We shall also require the following useful variant of the limit curve theorem. (See~\cite[Thm.\ 2]{EVP21} for a related result.) To this end, recall that a locally weakly causally closed, causally path-connected and localizable Lorentzian pre-length space is called a~\emph{Lorentzian length space\/} if
\begin{equation}
    \tau(x,y) = \sup \{L_\tau(\gamma) : \gamma \colon [a,b] \to X \text{ causal with }\gamma(a) = x, \, \gamma(b) = y\} \cup \{0\}.
\end{equation}
A Lorentzian length space is~\emph{globally hyperbolic\/} if it is non-totally imprisoning (i.e., $d$-lengths of causal curves in compact sets are uniformly bounded) and $J(x,y)$ is compact for all $x,y \in X$. Globally hyperbolic Lorentzian length spaces are strongly causal, cf.\ \cite[Thm.\ 3.26]{kunzinger2018lorentzian}.

\begin{Theorem}[Limit curve theorem: Converging endpoints]
\label{Theorem: Limitcurvetheoremtwoendpointcase}
Let $(X,d,\ll,\leq,\tau)$ be a globally hyperbolic Lorentzian length space. Let $\gamma_n \colon [0,1] \to X$ be a sequence of future-directed causal curves such that $p_n \coloneqq \gamma_n(0) \to p$ and $q_n \coloneqq \gamma_n(1) \to q$ with $p\ne q$. Then a subsequence of a reparametrized sequence $(\tilde{\gamma}_n)$ converges uniformly to a future-directed causal curve from $p$ to $q$. Moreover, we have:
\begin{enumerate}
\item\label{217i} If each $\gamma_n$ is a geodesic and $X$ is locally maximizing, then $\gamma$ is also a geodesic, and in this case $L_\tau(\gamma) = \lim_n L_\tau(\gamma_n)$. Moreover, if $X$ is regular and $p \ll q$, then $\gamma$ is a future-directed timelike geodesic.
\item\label{217ii} If $(\gamma_n)$ is a limit maximizing sequence (in particular: if each $\gamma_n$ is maximizing), then $\gamma$ is maximizing.
\end{enumerate}
\begin{proof}
    Choose $\Tilde{p} \in I^-(p)$ and $\Tilde{q} \in I^+(q)$. Then $\gamma_n([0,1]) \subseteq J(\Tilde{p}, \Tilde{q})$, which is compact. Since $X$ is non-totally imprisoning, there is a uniform upper bound on the $d$-arclengths of causal curves contained in $J(\Tilde{p},\Tilde{q})$, hence we may reparametrize the $\gamma_n$ with respect to $d$-arclength to get a uniformly Lipschitz sequence,
    (again denoted by the same letter) $\gamma_n \colon [a_n,b_n] \to X$, with both $(a_n)$ and $(b_n)$ bounded sequences. Without loss of generality we may suppose that $a_n\to a$ and $b_n\to b$. Let $\alpha_n$ be a future-directed causal curve connecting $\tilde{p}$ to $p_n=\gamma_n(a_n)$, parametrized by $d$-arclength, say $\alpha_n \colon [a_n-\delta_n,a_n] \to X$ for some $\delta_n>0$. Then $\delta_n \ge d(\tilde{p},p_n) \to d(\tilde{p},p)>0$. Similarly, let $\beta_n \colon [b_n,b_n+\eta_n] \to X$ be a $d$-arclength reparametrization of a future-directed causal curve from $q_n=\gamma_n(b_n)$ to $\tilde{q}$. Here, $\eta_n \ge d(q_n,\tilde{q}) \to d(q,\tilde{q})>0$. Hence we may suppose that $\min(\delta_n,\eta_n)>2c$ for some $c>0$ and all $n$. Let now $\hat{\gamma}_n$ denote the restriction of the concatenation of $\alpha_n$, $\gamma_n$ and $\beta_n$ to the interval $[a-c,b+c]$. Applying Theorem~\ref{Theorem: limitcurvetheoremforsegments} to this sequence it follows that we can extract a subsequence that converges uniformly to a causal curve whose $d$-arclength reparametrization is defined on $[a-c,b+c]$. 
    Denoting by $\gamma$ the restriction of this curve to $[a,b]$ we obtain a causal limit curve connecting $p$ to $q$.
   
    Part~\ref{217i} now follows from Proposition~\ref{Proposition: TopologyofGeoX} and the assumption that $p \neq q$. The additional claim is an immediate consequence of the fact that in regular spaces maximizers have a causal character.

    Part~\ref{217ii} follows from Theorem~\ref{Theorem: limitcurvetheoremforsegments}, which states that the curve $\gamma$ is maximizing if the sequence $\gamma_n$ is limit maximizing.
\end{proof}
\end{Theorem}

\begin{Remark}[Limit curve theorem for $\Geo(X)$]
\label{Remark: limitcurveconvendp for Geo}
We may reformulate Theorem~\ref{Theorem: Limitcurvetheoremtwoendpointcase} in the language of equivalence classes as follows: Let $X$ be a globally hyperbolic and locally maximizing Lorentzian length space and $\gamma_n \in \Geo(X)$. Suppose that $\gamma_n(0) \eqqcolon p_n \to p \neq q \leftarrow q_n \coloneqq \gamma_n(1)$. Then a subsequence of $(\gamma_n)$ converges in $\Geo(X)$ to some  $\gamma$. Moreover, if $X$ is regular and $p \ll q$, then $\gamma \in \TGeo(X)$. If each $\gamma_n$ is maximizing, then so is $\gamma$.
\end{Remark}

We end this subsection with a property of geodesic convergence in smooth semi-Riemannian manifolds. 
\begin{Lemma}\label{lem:geod_convergence}
Let $(M,g)$ be a smooth semi-Riemannian manifold. Let $a < 0 \leq 1 < b$. Then the following are equivalent for geodesic segments $\gamma_n, \gamma \colon [a,b] \to M$:
\begin{enumerate}
    \item $\gamma_n \to \gamma$ in $C^0$.
    \item $\gamma_n \to \gamma$ in $C^\infty$.
    \item $\gamma_n'(0) \to \gamma'(0)$ in $TM$.
\end{enumerate}
\end{Lemma}
\begin{proof}
Clearly, $(ii)$ implies $(i)$. To see that $(iii)$ implies $(ii)$, observe that if $(p_n,v_n) \in TM$ with $(p_n,v_n) \to (p,v)$, then the fact that $\gamma_n(t) \coloneqq \exp_{p_n}(t v_n)$ converges to $\gamma(t) \coloneqq \exp_{p}(tv)$ in $C^\infty_{\mathrm{loc}}$ on the maximal domain of $\gamma$ is a consequence of basic ODE theory, see e.g.\ \cite[Cor.\ 2.6]{KOSS}. We are left with showing that $(i)$ implies $(iii)$. Suppose $\gamma_n \colon [a,b] \to M$ converges uniformly to $\gamma \colon [a,b] \to M$. It is well-known that the exponential map $E \colon \mathcal{D} \subseteq TM \to M^2$ is a diffeomorphism from a neighborhood of the zero section in $TM$ onto a neighborhood of the diagonal in $M^2$. We may cover $\gamma$ with finitely many convex normal neighborhoods $U_i$ such that $\bigcup_i U_i^2$ is contained in the aforementioned neighborhood of the diagonal. For large enough $n$, the $U_i$ cover the corresponding portions of $\gamma_n$ due to uniform convergence. But then
    \begin{align*}
        \gamma_n'(0) = E^{-1}(\gamma_n(0),\exp_{\gamma_n(0)}(\gamma_n'(0))) = E^{-1}(\gamma_n(0),\gamma_n(1)) \to E^{-1}(\gamma(0),\gamma(1)) = \gamma'(0),
    \end{align*}
    concluding the proof.
\end{proof}


\section{Timelike conjugate points}
\label{Section: Timelikeconjugatepoints}

This section is dedicated to the discussion of several notions of conjugate points in the synthetic setting. These notions were introduced in positive definite signature by Shankar--Sormani~\cite{shankar2009conjugate} (cf.\ also Rinow~\cite{rinow1961innere}). We will for the most part follow their terminology and conventions when adapting the notions to the Lorentzian setting.

\subsection{Symmetric conjugate points}
\label{Subsection: symmetric conjugate points}

For the following discussion, let $(X,d,\ll,\leq, \tau)$ be a finitely $\tau$-measurable Lorentzian pre-length space.

\begin{Definition}[One-sided conjugate points]
Let $\gamma \in \TGeo(X)$ be a future-directed timelike geodesic from $p$ to $q$. We say $q$ is~\emph{one-sided conjugate to} $p$ if there is a sequence $q_i \to q$ and for each $i$, there exist two distinct future-directed timelike geodesics $\eta_i \neq \gamma_i \in \TGeo(X)$ from $p$ to $q_i$ such that $ \eta_i,\gamma_i \to \gamma$ with respect to $d_{\Gamma}$.
\end{Definition}

\begin{Theorem}\label{th:one-sided-vs-classical}
In a strongly causal smooth spacetime $(M,g)$, points $p\ll q$ are one-sided conjugate points along a future-directed timelike geodesic $\gamma$ if and only if they are conjugate points (in the Jacobi field sense).
\end{Theorem}

\begin{proof}
Let $p,q \in M$ be such that $p\ll q$, $\gamma$ a future-directed timelike geodesic from $p$ to $q$, and suppose that $q$ is one-sided conjugate to $p$ along $\gamma$.
Indirectly assuming that $p$ and $q$ are not conjugate along $\gamma$, then $\exp_p$ is regular around $v \coloneqq \exp_p^{-1}(q)$. In particular the exponential map is injective onto a neighborhood of $q$ and so there cannot exist points as given in the definition of one-sided conjugacy.

For the reverse implication we use a generalization, due to D.\ Szeghy~\cite{szeghy2008conjugate}, of a classical result by F.~W.\ Warner~\cite{warner1965conjugate}. Let $\gamma(t) = \exp_p(tv)$ be a timelike geodesic and let $q=\gamma(1)$ be conjugate to $p$ along
$\gamma$. Denoting by $C(p)\subseteq T_pM$ the conjugate locus of $p$ we thus have $v\in C(p)$. Let $U\subseteq T_pM$ be an open neighborhood of $v$ consisting entirely of timelike vectors. By 
\cite[Thm.\ 2.1]{szeghy2008conjugate}, together with~\cite[Rem.\ 4.1]{szeghy2008conjugate}, both the set $C^R(p)$ of
regular conjugate vectors and the set $C(p)'$ of conjugate vectors of locally minimal multiplicity
are open and dense in $C(p)$, hence so is their intersection $C^R(p)\cap C(p)'$.
Pick $v'\in C^R(p)\cap C(p)'\cap U$. Then by~\cite[Lem.\ 4.2]{szeghy2008conjugate}, $\exp_p$ is a Warner map on a suitable open neighborhood $V\subseteq U$ of $v'$ (note that by~\cite[Sec.\ 5]{szeghy2008conjugate}, the fact that $v'$ is timelike implies that $\mathcal{D}_{v'}$ is spacelike, so the assumptions of~\cite[Lem.\ 4.2]{szeghy2008conjugate} are indeed satisfied). Therefore (cf.\ \cite[Thm.\ 2.2]{szeghy2008conjugate} and the remarks following it), $\exp_p$ is not injective on any neighborhood of $v'$. In particular, it is not injective on $U$.
Consequently, there exist vectors $v_i, w_i \in U, v_i\neq w_i$ with $v_i, w_i\to v$ and $\exp_p(v_i)=\exp_p(w_i) \eqqcolon q_i$. The distinct timelike geodesics $\gamma_i(t) \coloneqq \exp_p(t v_i)$ and $\sigma_i(t) \coloneqq \exp_p(t w_i)$ then satisfy one-sided conjugacy.
\end{proof}

Next, we discuss a notion of conjugate point similar to one-sided conjugacy, but here the initial point is also allowed to vary.

\begin{Definition}[Symmetric conjugate points]
Let $p\ll q$. Then $p,q$ are called~\emph{symmetrically conjugate} along a future-directed timelike geodesic $\gamma \in \TGeo(X)$ running between them if there are $p_i \to p$ and $q_i \to q$ such that for each $i$ there are two distinct future-directed timelike geodesics $\gamma_i \neq \sigma_i \in \TGeo(X)$ from $p_i$ to $q_i$ such that $\gamma_i \to \gamma$ and $\sigma_i \to \gamma$ with respect to $d_{\Gamma}$.
\end{Definition}
Clearly, any one-sided conjugate point is also symmetrically conjugate. The converse is open even in the metric case, cf.\ \cite[Open Problem 2.5]{shankar2009conjugate}.

\begin{Remark}[Symmetric conjugate points - trivial case]
Let $X$ be a finitely $\tau$-measurable
Lorentzian pre-length space. If $\gamma, \sigma \in \TGeo(X)$ are different geodesics with the same endpoints $p \coloneqq \sigma(0) = \gamma(0)$ and $q \coloneqq \sigma(1) = \gamma(1)$, then $p,q$ are one-sided conjugate (and hence also symmetric conjugate).
\end{Remark}


\begin{Proposition}[Equivalence of notions]\label{prop:equivalence_of_conjugate_notions}
Let $(M,g)$ be a strongly causal smooth spacetime and let $\gamma$ be a future-directed timelike geodesic from $p$ to $q$. Then the following are equivalent:
\begin{enumerate}
    \item[(i)] $q$ is one-sided conjugate to $p$ along $\gamma$.
    \item[(ii)] $p,q$ are symmetrically conjugate along $\gamma$.
    \item[(iii)] $q$ is a conjugate point of $p$ along $\gamma$ in the classical sense.
\end{enumerate}
\end{Proposition}
\begin{proof}
By Theorem~\ref{th:one-sided-vs-classical}, (i)$\Leftrightarrow$(iii). As observed above, one-sided conjugate implies symmetric conjugate, i.e.\ (i)$\Rightarrow$(ii). Finally, (ii)$\Leftrightarrow$(iii) follows by the same 
proof as in~\cite[Thm.\ 2.6]{shankar2009conjugate}.
\end{proof}

\subsection{Families of timelike geodesics}
\label{subsection: Families of timelike geodesics}

\begin{Definition}[Continuous timelike families]
Let $\gamma\in \TGeo(X)$ be a future-directed timelike geodesic from $p=\gamma(0)$ to $q=\gamma(1)$. A map $F \colon U\times V\to \TGeo(X)$
with $U$ a neighborhood of $p$ and $V$ a neighborhood of $q$ is called a~\emph{timelike family about $\gamma$\/} if 
$U\ll V$ (i.e., $p' \ll q'$ for all $p' \in U$ and $q' \in V$) and $F(p',q')$ goes from $p'$ to $q'$ for each $(p',q')\in U\times V$
and $F(p,q)=\gamma$. 
We say that $F$ is~\emph{continuous about\/} $\gamma$ if $F$ is a continuous (w.r.t.\ $d_{\Gamma}$) map.
$F$ is called~\emph{continuous at\/} $\gamma$ if it is only required to be continuous (w.r.t.\ $d_{\Gamma}$) at the point $(p,q)$ itself.
\end{Definition}
Observe that since $\ll$ is an open relation (cf.~\cite[Prop.\ 2.13]{kunzinger2018lorentzian}), whenever $x\ll y$ there exist neighborhoods $U$ of $x$ and $V$ of $y$ such that $p' \ll q'$ for all $p' \in U$ and $q' \in V$.

\begin{Definition}[Unique timelike families]\label{def:unique_familiy}
Let $\gamma$ be a future-directed timelike geodesic. We say that there exists a unique timelike
family which is continuous about $\gamma$ resp.\ at $\gamma$ if
\begin{itemize}
\item[(i)] there exists a timelike family $F \colon U\times V \to \TGeo(X)$ that is continuous about $\gamma$ resp.\ at $\gamma$, and
\item[(ii)] if $F' \colon U' \times V' \to \TGeo(X)$ is another timelike family about $\gamma$ that is continuous about $\gamma$ resp.\ at $\gamma$, then there exist neighborhoods $U''\subseteq U\cap U'$ and $V''\subseteq V\cap V'$ of $\gamma(0)$
resp.\ $\gamma(1)$ such that $F$ coincides with $F'$ on $U''\times V''$.
\end{itemize}
\end{Definition}

Another way to express this is to say that there exists a unique germ of timelike
families around $(\gamma(0), \gamma(1))$ with the respective continuity property. Given a timelike family $F \colon U\times V \to \TGeo(X)$ about $\gamma$, we shall denote by $[F]$ the germ around $(\gamma(0), \gamma(1))$ determined by $F$. Note that if a timelike family is continuous about $\gamma$ and unique among timelike families that are continuous at $\gamma$, then it is also unique among timelike families that are continuous about $\gamma$.

A $C^0$-continuous family is a family as above that is continuous w.r.t.\ $d_F$ instead of $d_{\Gamma}$. As we saw before, if $X$ is localizable, locally weakly causally closed, strongly causal and locally maximizing, then these two notions coincide, cf.\ Proposition~\ref{Proposition: TopologyofGeoX}. 

The following basic observation will repeatedly be used below.

\begin{Lemma}\label{lem:no-unique-fam-conj-points} Let $\gamma\in \TGeo(X)$ connect the points $p$ and $q$ and suppose that there exists a timelike family $F \colon U\times V \to \TGeo(X)$ continuous at $\gamma$. If there isn't a unique timelike family continuous at $\gamma$, then $p$ and $q$ are symmetrically conjugate. 
\end{Lemma}
\begin{proof} Our assumptions mean that (the second alternative in) (i) in Definition~\ref{def:unique_familiy} is satisfied, but (ii) is not. Hence there exists some timelike family $F' \colon U'\times V' \to \TGeo(X)$ continuous at $\gamma$ that does not coincide with $F$ on sufficiently small neighborhoods of $(p,q)$. For each $i\in \N$ pick neighborhoods $U_i\subseteq U\cap U'$ of $p$ and $V_i\subseteq V\cap V'$ such that $F|_{U_i\times V_i} \ne F'|_{U_i\times V_i}$, as well as
$U_i\subseteq B^d_{1/i}(p)$ and $V_i\subseteq B^d_{1/i}(q)$. Furthermore, by continuity of $F$ and $F'$ at $(p,q)$ we may assume 
$U_i$ and $V_i$ small enough to grant that $F(U_i\times V_i)
\subseteq B^{d_\Gamma}_{1/i}(\gamma)$ and also $F'(U_i\times V_i)
\subseteq B^{d_\Gamma}_{1/i}(\gamma)$. Pick $(p_i,q_i) \in U_i\times V_i$ such that
$\gamma_i \coloneqq F(p_i,q_i)\ne F'(p_i,q_i) \eqqcolon \sigma_i$. Then $\gamma_i$ and $\sigma_i$ are timelike
geodesics as required to establish that $p$ and $q$ are symmetrically conjugate.
\end{proof}


As an immediate consequence we obtain the following analogue of~\cite[Thm.\ 3.7]{shankar2009conjugate}
\begin{Corollary}\label{cor:not-sym-con-unique-germ}
If $p, q\in X$ are not symmetrically conjugate along $\gamma\in \TGeo(X)$ and if there exists a timelike family $F \colon U\times V \to \TGeo(X)$ continuous at $\gamma$, then the germ of $F$ at $(p,q)$ is the unique germ of a timelike family continuous at $\gamma$. 
\end{Corollary}

Collecting the previous results, we obtain:
\begin{Proposition}\label{prop:symm-conj-characterization} Let $\gamma\in \TGeo(X)$ connect $p$ and $q$ and suppose there exists a timelike family $F \colon U\times V \to \TGeo(X)$ continuous at $\gamma$. Then the following are equivalent:
\begin{itemize}
\item[(i)] $p$ and $q$ are not symmetrically conjugate along $\gamma$.
\item[(ii)] $F$ is a representative of the unique germ of timelike families continuous at $\gamma$.
\end{itemize}
\end{Proposition}

\begin{Definition}[Embeddability] A sequence $(\gamma_i)_{i\in \N}$ in $\TGeo(X)$ connecting points $p_i$ and $q_i$ and converging (with respect to $d_\Gamma$) to a curve $\gamma\in \TGeo(X)$ is~\emph{embeddable\/} if there exists a timelike family $F \colon U\times V \to \TGeo(X)$ which is continuous at $\gamma$ such that, for $i$ large, $(p_i,q_i)\in U\times V$ and $F(p_i,q_i) = \gamma_i$. It is~\emph{strongly embeddable\/} if the timelike family can be chosen to be continuous about $\gamma$.
\end{Definition}
In particular a constant sequence $(\gamma)$ is embeddable resp.\ strongly embeddable if and only if there exists a timelike family about $\gamma$ that is continuous at $\gamma$ resp.\ about $\gamma$.
In this case we say that $\gamma$ itself is (strongly) embeddable. 

\begin{Proposition}\label{prop:embed-char}
Let $\gamma\in \TGeo(X)$ connect $p$ to $q$. Then the following are equivalent:
\begin{itemize}
\item[(i)] $\gamma$ is embeddable.
\item[(ii)] For all sequences $p_i\to p$ and $q_i\to q$ there exists $i_0\in \N$ such that for all $i\ge i_0$ there is a geodesic $\gamma_i\in \TGeo(X)$ from $p_i$ to $q_i$ such that 
$\gamma_i\to \gamma$ with respect to $d_\Gamma$.
\end{itemize}
\end{Proposition}
\begin{proof}
(i)$\Rightarrow$(ii): Let $F \colon U \times V$ be a timelike family continuous at $\gamma$. Then given sequences as in (ii), for $i$ large
we may set $\gamma_i \coloneqq F(p_i,q_i)$. \\
(ii)$\Rightarrow$(i): We first show that for each $k\in \N$ there exist neighborhoods $U_k$ of $p$ and $V_k$ of $q$ such that for all $(p',q')\in U_k\times V_k$ there exists a timelike geodesic $\sigma\in \TGeo(X)$ from $p'$ to $q'$ with $d_\Gamma(\gamma,\sigma)<\frac{1}{k}$. Indeed, if this were not the case, then for some $k\in \N$ there would exist a sequence $(p_i,q_i)\to (p,q)$ such that
any timelike geodesic connecting $p_i$ to $q_i$ would have $d_\Gamma$-distance from $\gamma$ greater than $\frac{1}{k}$, contradicting (ii).
Since the relation $\ll$ is open, we can arrange that $U_k\ll V_k$. Furthermore, we may assume that $U_{k+1}\subseteq U_k$ and $V_{k+1}\subseteq V_k$ for all $k$ and that $\bigcap_k U_k = \{p\}$, as well as $\bigcap_k V_k = \{q\}$. Now set $W_k \coloneqq (U_k\times V_k)\setminus (U_{k+1}\times V_{k+1})$ for each $k$ and define 
\[
F \colon U_1 \times V_1 \to \TGeo(X)
\]
as follows: Set $F(p,q) \coloneqq \gamma$, and for each $(p',q')\in (U_1\times V_1)\setminus \{(p,q)\}$ take the unique
$k\in \N$ such that $(p',q')\in W_k$ and define $F(p',q')$ to be any timelike geodesic from $p'$ to $q'$ satisfying $d_\Gamma(\gamma,F(p',q'))<\frac{1}{k}$. Then $F$ is a timelike family, and it is also continuous at $\gamma$: if $(p_i,q_i)\to (p,q)$, then for any $k\in \N$ there exists some $i_0$
such that $(p_i,q_i)\in U_k\times V_k$ for all $i\ge i_0$, implying that $d_\Gamma(\gamma,F(p_i,q_i))<\frac{1}{k}$.
\end{proof}

\begin{Lemma} Suppose that every $d_\Gamma$-convergent sequence $(\gamma_i)_{i\in \N}$ in $\TGeo(X)$ is embeddable.
If for every future-directed timelike geodesic $\gamma \in \TGeo(X)$ such that $L_{\tau}(\gamma) < R$ there exists a unique timelike family about $\gamma$ that is continuous at $\gamma$, then it is even true that for every future-directed timelike geodesic $\gamma$ with $L_{\tau}(\gamma) < R$ there exists a timelike family that is continuous about $\gamma$.
\end{Lemma}
\begin{proof}
Let $F \colon U \times V \to\TGeo(X)$ be a timelike family about a future-directed timelike geodesic $\gamma \colon [0,1] \to X$ that is continuous at $(\gamma(0),\gamma(1))$. By $d_{\Gamma}$-continuity of $F$ at $(\gamma(0),\gamma(1))$, we can shrink $U$ and $V$ to assume w.l.o.g.\ that all geodesics $F(x,y)$ satisfy $L_{\tau}(F(x,y)) < R$, $x \in U$, $y \in V$.\\
Suppose there are no neighborhoods $U',V'$ of $\gamma(0),\gamma(1)$ (with $U'\ll V'$) such that $F|_{U' \times V'}$ is continuous. Then there are sequences $p_i \to \gamma(0)$ and $q_i \to \gamma(1)$ such that $F$ is not continuous at any $(p_i,q_i)$. Let $\gamma_i \coloneqq F(p_i,q_i)$. So for any $i$ there are $\varepsilon_i>0$ and sequences $p_{ij} \to p_i$ and $q_{ij} \to q_i$ such that
\begin{align}\label{eq:nolim}
    d_\Gamma(F(p_{ij}, q_{ij}),F(p_i,q_i))>\varepsilon_i.
\end{align}
Since $L_{\tau}(\gamma_i) < R$, for each $\gamma_i$ there is a timelike family $F_i$ about $\gamma_i$ continuous at its endpoints $(p_i,q_i)$. In particular,
\begin{align}
    \lim_j F_i(p_{ij},q_{ij}) = F_i(p_i,q_i)=\gamma_i = F(p_i,q_i).
\end{align}
So by \eqref{eq:nolim} for each $i$ there is $j_i$ such that for all $j \geq j_i$,
\begin{align}
    d_\Gamma(F(p_{ij},q_{ij}), F_i(p_{ij},q_{ij})) > \varepsilon_i/2.
\end{align}
We can choose $j_i$ larger to guarantee that 
\begin{align}
    d_{\Gamma}(F_i(p_{ij_i},q_{ij_i}),\gamma_i) < 1/i
\end{align}
due to the convergence noted above, and also $(p_{ij_i},q_{ij_i}) \to (\gamma(0),\gamma(1))$ by a diagonal argument. Thus we have two distinct sequences of future-directed timelike geodesics, namely $F(p_{ij_i},q_{ij_i})$ and $F_i(p_{ij_i},q_{ij_i})$ both converging to $\gamma$, where for the second convergence we use
\begin{align}
 d_\Gamma(F_i(p_{ij_i}, q_{i,j_i}),\gamma)\leq 1/i+d_\Gamma(F(p_i,q_i),F(p,q)).
\end{align}
Now embed the sequence $(F_i(p_{ij_i}, q_{i,j_i}))_i$ of geodesics into a timelike family $F'$ that is continuous at $\gamma$. Then $[F]\ne [F']$, contradicting our assumption. 
\end{proof}

\begin{Remark} The previous result is a direct analogue of~\cite[Lem.\ 3.4]{shankar2009conjugate}. Note, however, that the embeddability assumption we make is absent from that result, although we are under the impression that it is needed to finish the proof.  
\end{Remark}

\begin{Proposition}\label{prop:embeddable_for_spacetimes} Let $(M,g)$ be a strongly causal smooth spacetime. Let $(\gamma_i)_{i\in\mathbb{N}}$ be a sequence of timelike geodesics in $M$ converging (with respect to $d_{\Gamma}$) to a limit geodesic $\gamma$. Let $p = \gamma(0)$ and $q = \gamma(1)$. 
If $p$ and $q$ are not conjugate along $\gamma$, then the sequence $(\gamma_i)$ is strongly embeddable.
\end{Proposition}
\begin{proof} 
Let $v_i = \gamma_i'(0) \in TM$ and $v = \gamma'(0) \in T_pM$. By Lemma \ref{lem:geod_convergence},
we have $v_i \to v$. Consider the map $E: TM \supseteq \mathcal{D} \to M\times M$ defined on
a suitable open neighbourhood $\mathcal{D}$ of the $0$-section in $TM$ by 
\[ E(w) = (\pi(w), \exp_{\pi(w)}(w)), \]
with $\pi: TM \to M$ the bundle projection. Then the differential of $E$ at $v$ takes the following lower-triangular block form:
\[
    T_vE \cong \begin{pmatrix}
    \mathrm{id} & 0 \\
    \star & T_v\exp_p
    \end{pmatrix}.
\]
As $p$ and $q$ are not conjugate along $\gamma$, $T_v\exp_p$ is invertible, implying that $E$ is a diffeomorphism from a neighbourhood $W \subset TM$ of $v$ onto a neighbourhood $U \times V$ of $(p, q)$ in $M\times M$. Let $H: U \times V \to W$ be the smooth inverse of $E|_{W}$.

We define the timelike family $F: U \times V \to \TGeo(M)$ by
\[ F(x, y)(t) := \exp_x(t \cdot H(x, y)). \]
This family is continuous about $\gamma$. Moreover, since $\gamma_i \to \gamma$, for sufficiently large $i$ the initial velocities $v_i$ lie in $W$ and the endpoints $(p_i, q_i)$ lie in $U \times V$. Consequently, $H(p_i, q_i) = v_i$, and $F(p_i, q_i)$ is exactly the geodesic $\gamma_i$. Thus, the sequence is strongly embeddable.
\end{proof}

The following result is a direct analogue of \cite[Prop.\ 3.5]{shankar2009conjugate}.
\begin{Theorem}\label{thm:classical-conj-vs-unique-family}
Let $(M,g)$ be a smooth strongly causal spacetime and let $\gamma \colon [0, 1] \to M$ be a future-directed timelike geodesic connecting $p$ to $q$. Then the following are equivalent:
\begin{itemize}
\item[(i)] $p$ and $q$ are not conjugate (in the classical sense) along $\gamma$.
\item[(ii)] There exists a timelike family which is continuous about $\gamma$ and which is
unique among all timelike families that are continuous at $\gamma$.
\end{itemize}
\end{Theorem} 
\begin{proof}
(i)$\Rightarrow$(ii): Applying Proposition~\ref{prop:embeddable_for_spacetimes} to the constant sequence $(\gamma)$
it follows immediately from the proof of that result that the timelike
family constructed there is continuous about $\gamma$ and is unique among all such families.

(ii)$\Rightarrow$(i): Proposition~\ref{prop:equivalence_of_conjugate_notions} shows that (i) is equivalent to $p$ and $q$ not
being symmetrically conjugate. The claim therefore follows from Proposition~\ref{prop:symm-conj-characterization}.
\end{proof}

\begin{Proposition}[Existence of families around unique maximizers]
\label{Proposition: Existenceoffamiliesifgeodunique}
Let $(X,d,\ll,\leq,\tau)$ be a regular, globally hyperbolic and locally maximizing Lorentzian length space. Let $p \ll q$ and suppose there is a unique maximizer $\gamma \in \TGeo(X)$ from $p$ to $q$. Then there exists a timelike family about $\gamma$ consisting of maximizers which is continuous at $\gamma$, in particular, $\gamma$ is embeddable.
\begin{proof}
    Let $U, V$ be neighborhoods of $p,q$ so that $x \ll y$ for all $x \in U$, $y \in V$. For each $(x,y) \in U \times V$, let $F(x,y)$ be a maximizer from $x$ to $y$, hence $F(p,q)=\gamma$. We claim that $F$ is continuous at $\gamma$. Indeed, if $(p_i,q_i)\to (p,q)$, then by Theorem~\ref{Theorem: Limitcurvetheoremtwoendpointcase} (see also Remark~\ref{Remark: limitcurveconvendp for Geo}) and regularity of $X$, 
    there exists a subsequence of $(F(p_i,q_i))$ that converges (in $d_{\Gamma}$) to a timelike maximizer connecting $p$ to $q$. By uniqueness, this maximizer has to coincide with $\gamma$. Since the same is true for any subsequence of $(F(p_i,q_i))$, we must have that $d_\Gamma(F(p_i,q_i),F(p,q))\to 0$ as $i\to \infty$.  
\end{proof}
\end{Proposition}


\subsection{Unreachable and ultimate conjugate points}
\label{Subsection: unreachable and ultimate}

We continue to assume that $X$ is a finitely $\tau$-measurable Lorentzian pre-length space.

\begin{Definition}[Unreachable conjugate points]
Let $\gamma \in \TGeo(X)$ with $\gamma(0) = p$ and $\gamma(1) = q$. We say that $p,q$ are~\emph{unreachable conjugate along\/} $\gamma$ if $\gamma$ is not embeddable.
Otherwise, $p$ and $q$ are called~\emph{reachable\/} along $\gamma$.
\end{Definition}


Proposition~\ref{prop:embed-char} immediately yields the following characterization:
\begin{Proposition}\label{prop:unreachable-characterization} 
Let $\gamma\in \TGeo(X)$ connect $p$ and $q$. Then the following are equivalent:
\begin{itemize}
\item[(i)] $p$ and $q$ are unreachable conjugate along $\gamma$; 
\item[(ii)] There exist sequences $p_i\to p$ and $q_i\to q$ such that for no choice $\gamma_i$ of timelike
geodesic from $p_i$ to $q_i$ we have that $\gamma_i \to \gamma$ with respect to $d_\Gamma$.
\end{itemize}
\end{Proposition}
\begin{Lemma}\label{lem:unreachable_implies_standard_in_spacetimes}
Let $(M,g)$ be a strongly causal
spacetime. If $p\ll q$ are unreachable conjugate along a timelike geodesic $\gamma$, then they are conjugate (in the classical sense) along $\gamma$. 
\end{Lemma}
\begin{proof}
If $p$ and $q$ were not conjugate in the classical sense, then
by Theorem~\ref{thm:classical-conj-vs-unique-family} there would exist a
timelike family about $\gamma$ that is continuous at $\gamma$, contradicting
Proposition~\ref{prop:unreachable-characterization} (iii).
\end{proof}

\begin{Example}
    Consider $n$-dimensional anti-de Sitter space $H^{1,n}(K)$ of conjugate radius $\frac{\pi}{\sqrt{-K}}$, $K < 0$. Then any two timelike related points $p \ll q$ with $\tau(p,q) = \frac{\pi}{\sqrt{-K}}$ are unreachable conjugate along any timelike maximizer between them. Indeed, any future-directed timelike geodesic $\gamma$ containing $p$ and $q$ lies in the intersection of $H^{1,n}(K)$ with some plane $P$ connecting $p$, $q$ and $0$. Now let $Q\ne P$ be another plane through $0$, $p$ and $q$ such that $Q$ cuts out from $H^{1,n}(K)$ a future-directed timelike geodesic $\sigma$ from $p$ to $q$.  Then for points $p_i, q_i \in \sigma$ with $p\ll p_i$, $q_i\ll q$, $p_i\to p$ and $q_i\to q$, $\sigma$ is the only future-directed timelike geodesic connecting them.
\end{Example}

\begin{Corollary}\label{cor:unique-max-reachable}
    Let $(X,d,\ll,\leq,\tau)$ be a regular, globally hyperbolic and locally maximizing Lorentzian length space. Let $p \ll q$ and suppose there exists a unique maximizer $\gamma \in \TGeo(X)$ from $p$ to $q$. Then $p,q$ are reachable.
\begin{proof}
    This follows from Proposition~\ref{Proposition: Existenceoffamiliesifgeodunique}.
\end{proof}
\end{Corollary}

\begin{Definition}[Ultimate conjugate points]
Let $(X,d,\ll,\leq,\tau)$ be a finitely $\tau$-measurable Lorentzian pre-length space and $\gamma \in \TGeo(X)$. Then $\gamma(0),\gamma(1)$ are called~\emph{ultimate conjugate points\/} if they are unreachable conjugate points or symmetric conjugate points (along $\gamma$).
\end{Definition}

\begin{Example} In~\cite[Ex.\ 4.7]{shankar2009conjugate}, an example (due to V.\ Schroeder) of a geodesic $\sigma \colon [0,1] \to N$ in a
complete Riemannian manifold $(N,h)$ (specifically, a two-dimensional ellipsoid) is given, with the following properties: $\sigma$ is the unique geodesic in $N$ connecting $x=\sigma(0)$ and $y=\sigma(1)$, while $x$ and $y$ are conjugate along $\sigma$. We may ``Lorentzify'' this example as follows: Let $M \coloneqq \R\times N$ with Lorentzian metric $g=-dt^2 + h_x$. Then $(M,g)$ is globally hyperbolic (each slice $\{t=t_0\}$ is a Cauchy surface), and the following correspondences hold (cf., e.g., \cite[Ch.\ 7]{ON83}): Any geodesic in $(M,g)$ is of the form $s\mapsto (as+b,\beta)$, where $a,b\in \R$ and $\beta$ is a geodesic in $(N,h)$. Jacobi fields in $(M,g)$ take the form $s\mapsto (T(s),J(s))$, where $T$ is affine linear and $J$ is a Jacobi field in $(N,h)$. Thus if we set $\gamma(s) \coloneqq (as,\sigma(s))$ with $\sigma$ as above, then for $a>0$ sufficiently large, $\gamma$ is a timelike geodesic connecting $p=(0,x)$ to $q=(a,y)$. It is the unique timelike geodesic in $(M,g)$ connecting $p$ and $q$, and these points are conjugate along $\gamma$. Since $\sigma$ is minimizing from $x$ to $y$ in $N$, $\gamma$ is a timelike maximizer from $p$ to $q$ in $(M,g)$. By Corollary~\ref{cor:unique-max-reachable}, $p$ and $q$ are reachable along $\gamma$. Consequently, they provide an example of points along a timelike geodesic that are (by Proposition~\ref{prop:equivalence_of_conjugate_notions}) symmetrically (hence ultimate) conjugate, but not unreachable conjugate.
\end{Example}

\begin{Corollary}
Let $(M,g)$ be a smooth, strongly causal spacetime and $\gamma \colon [0,1] \to M$ a timelike geodesic. Then $\gamma(0),\gamma(1)$ are conjugate in the sense of Jacobi fields if and only if they are ultimate conjugate.
\begin{proof} This is immediate from Proposition~\ref{prop:equivalence_of_conjugate_notions} and Lemma 
\ref{lem:unreachable_implies_standard_in_spacetimes}.
\end{proof}
\end{Corollary}

\begin{Proposition}
    Let $(X,d,\ll,\leq,\tau)$ be a regular, locally maximizing, globally hyperbolic Lorentzian length space. Let $\gamma \in \TGeo(X)$ such that $p \coloneqq \gamma(0)$ and $q \coloneqq \gamma(1)$ are not ultimate conjugate along $\gamma$. Then there exists a continuous timelike family around $\gamma$ whose germ is unique among all timelike families continuous at $\gamma$.
    \begin{proof}
        By assumption, $p,q$ are reachable and not symmetric conjugate. By definition of reachability, there exists a timelike family $F \colon U \times V \to \TGeo(X)$ continuous at $\gamma$. By Proposition~\ref{prop:symm-conj-characterization}, $F$ is a representative of the unique germ of timelike families continuous at $\gamma$.
    \end{proof}
\end{Proposition}


\subsection{Causal cut loci}
\label{Subsection: causalcutloci}

\begin{Definition}[Causal cut points]
Let $(X,d,\ll,\leq,\tau)$ be a finitely $\tau$-measurable Lorentzian pre-length space and $p,q \in X$. We say $q$ is a~\emph{future causal cut point\/} of $p$ if there exist at least two distinct maximizing elements in $\Geo(X)$ connecting $p$ to $q$. Moreover, if $X$ is regular, we call a future causal cut point~\emph{timelike\/} if $p \ll q$ and~\emph{null\/} if $p < q$.
\end{Definition}

Recall that geodesics in regular Lorentzian pre-length spaces are either timelike or null by definition. Hence, the definition of future causal cut points implies that in such spaces, timelike cut points are joined by two distinct timelike maximizers, and similarly null cut points are joined by two distinct null maximizers.

\begin{Definition}[Initial cut loci]
\label{Definition: firstcutlocus}
Let $(X,d,\ll,\leq,\tau)$ be a finitely $\tau$-measurable Lorentzian pre-length space and $p \in X$. Given $\gamma \in \Geo(X)$, we say $q \coloneqq \gamma(1)$ is an~\emph{initial future causal cut point\/} of $p=\gamma(0)$ along $\gamma$ if
\begin{enumerate}
    \item there exist future causal cut points $q_i$ of $p$ converging to $q$;
    \item there is no point $q' \in \gamma([0,1]) \setminus \{q\}$ with property (i).
\end{enumerate}
In the case where there is always a $q'$ as in (ii), we take $p$ to be the initial future causal cut point of itself along $\gamma$. We denote the set of initial future causal cut points of $p$ along causal geodesics by $\mathrm{IniCut}_C^+(p) \subseteq X$, called the~\emph{initial future causal cut locus\/} of $p$.
If $X$ is regular, we define the~\emph{initial future timelike cut locus} $\mathrm{IniCut}_T^+(p)$ and the~\emph{initial future null cut locus} $\mathrm{IniCut}_N^+(p)$ by replacing all instances of causal with timelike resp.\ null. It is then clear that $\mathrm{IniCut}_C^+(p)$ is the disjoint union of $\mathrm{IniCut}_T^+(p)$ and $\mathrm{IniCut}_N^+(p)$.
\end{Definition}

\begin{Remark}[Cut points in smooth spacetimes]
Let $(M,g)$ be a smooth, globally hyperbolic spacetime. Recall that if $\gamma \colon [0,1] \to M$ is a future causal geodesic, then $\gamma(1)$ is in the future causal cut locus (in the standard sense) of $p$ if and only if there exists another maximizer from $p$ to $q$ ($\gamma$ is also maximizing up to $q$) or $q$ is the first conjugate point to $p$ along $\gamma$ (the former case corresponds $q_i = q$ in Definition~\ref{Definition: firstcutlocus}, see e.g.\ \cite[Prop.\ 3.2.30]{Treude:2011}). 
Hence, $\mathrm{IniCut}_C^+(p)$ agrees with the usual notion of future causal cut locus in this setting. Similarly, $\mathrm{IniCut}_T^+(p)$ and $\mathrm{IniCut}_N^+(p)$ agree with the usual notions of future timelike and future null cut locus, respectively.
\end{Remark}

\begin{Remark}[Comparison with~\cite{shankar2009conjugate}]
In~\cite[Def.\ 7.2]{shankar2009conjugate}, the authors require the $q_i$ from Definition~\ref{Definition: firstcutlocus} (in their metric setting) to lie on $\gamma$. 
This notion, however, does not coincide with the standard notion on smooth Riemannian manifolds.
Indeed, the two-dimensional torus of revolution provides a counterexample: The cut locus of a point $p$ slightly above the outer equator contains an arc lying on a parallel whose endpoints are conjugate but cannot be reached by more than one minimizer from $p$, see \cite{GMST:05}, (iii) on p.\ 104 and Fig.\ 4. Consider one of these (two) points, call it $q$ and suppose it is reached from $p$ by a geodesic $\gamma$. Since $\gamma$ beyond $q$ cannot coincide with the arc it does not contain any meeting points of minimizers from $p$ (although $q$ is a limit of such points lying on the arc).
\end{Remark}

\begin{Definition}[Initial injectivity radii]
Let $(X,d,\ll,\leq,\tau)$ be a finitely $\tau$-measurable Lorentzian pre-length space. For $p \in X$, we define the~\emph{initial future causal injectivity radius} of $X$ at $p$ to be $\mathrm{IniInj}_C^+(p) \coloneqq d(p,\mathrm{IniCut}_C^+(p))$. The~\emph{initial future causal injectivity radius} of $X$ is then $\mathrm{IniInj}_C^+(X) \coloneqq \inf_{p \in X} \mathrm{IniInj}_C^+(p)$. If $X$ is in addition regular, we define in complete analogy the~\emph{initial future timelike/null injectivity radii} $\mathrm{IniInj}_T^+(p),\mathrm{IniInj}_T^+(X), \mathrm{IniInj}_N^+(p)$ and $\mathrm{IniInj}_N^+(X)$.
\end{Definition}

\begin{Definition}[Unique injectivity radii]
Let $(X,d,\ll,\leq,\tau)$ be a finitely $\tau$-measurable and geodesic Lorentzian pre-length space. Then the~\emph{future causal unique injectivity radius} $\mathrm{UniqueInj}_C^+(X)$ is the supremum over all $r \geq 0$ so that all causally related points of $d$-distance less than $r$ can be connected by unique maximizers. If $X$ is in addition regular, we define and denote the corresponding timelike and null notions in an analogous way.
\end{Definition}

Our main result of this subsection is the following compatibility property.

\begin{Proposition}
Let $(X,d,\ll,\leq,\tau)$ be a finitely $\tau$-measurable and geodesic Lorentzian pre-length space. Then $\mathrm{IniInj}_C^+(X) = \mathrm{UniqueInj}_C^+(X)$ (timelike, null as well).
\end{Proposition}
\begin{proof}
Let $r \geq 0$ be such that all causally related points of distance at most $r$ can be connected by a unique maximizer, let $p \in X$ and $q \in
\mathrm{IniCut}_C^+(p)$.
Then by definition, either there exist two maximizers from $p$ to $q$ (in which case $d(p,q) \geq r$), or, arbitrarily close to $q$ there exist cut points of $p$. In the latter case, if $d(p,q) < r$, let $\varepsilon > 0$ and $\Tilde{q}$ a cut point of $p$ so that $d(p,\Tilde{q}) \leq d(p,q) + \varepsilon < r$, which is a contradiction. So in any case, $d(p,q) \geq r$, and taking the infimum over $q$ and then $p$, as well as the supremum over $r \geq 0$, we get
\begin{align}\label{eq:first-ultimate}
    \mathrm{IniInj}_C^+(X)
    \geq \mathrm{UniqueInj}_C^+(X).
\end{align}
Conversely, let $R \coloneqq \mathrm{UniqueInj}_C^+(X)$ and let $\varepsilon > 0$ be arbitrary. By definition, there exists $p \in X$ and $q \in J^+(p) \cap B_{R+\varepsilon}(p)$ so that there exist two distinct maximizers from $p$ to $q$. Therefore,
\[ 
\mathrm{IniInj}_C^+(X)
\le 
\mathrm{IniInj}_C^+(p)
\le d(p,q) < R+\eps.
\]
Letting $\eps\to 0$ shows that the converse of \eqref{eq:first-ultimate}.
\end{proof}

\section{Applications}
\label{Section: applications}

In this section, we give two applications that use of our notions of conjugate points, namely a timelike Rauch comparison theorem (in analogy with Shankar--Sormani~\cite[Thm.\ 10.1]{shankar2009conjugate}), as well as a consequence of the Lorentzian Cartan--Hadamard Theorem established by Er\"{o}s--Gieger~\cite{Eroes-Gieger25} (for the metric Cartan--Hadamard Theorem, cf.\ Bridson--Haefliger~\cite[Thm.\ 4.1]{bridson2013metric}).

\subsection{A timelike Rauch comparison theorem}
\label{subsection: Timelike Rauch}

\begin{Theorem}[Timelike Rauch comparison theorem]
\label{Theorem: Timelike Rauch}
Let $(X,d,\ll,\leq,\tau)$ be a chronological, regular and locally timelike maximizing Lorentzian pre-length space with timelike curvature $\leq \kappa$ and let $\gamma \in \TGeo(X)$ be such that $p \coloneqq \gamma(0)$ and $q \coloneqq \gamma(1)$ are symmetric conjugate along $\gamma$. Then $L_{\tau}(\gamma) \geq D_{\kappa}$.
\begin{proof}


    Suppose that there exists a timelike geodesic $\gamma \in \TGeo(X)$ with $L_{\tau}(\gamma) < D_{\kappa}$ whose endpoints are symmetric conjugate. Pick a representative, i.e., fix a parametrization $\gamma \colon [0,1] \to X$.

Cover $\gamma([0,1])$ by a finite collection $\mathcal{U}$ of locally timelike maximizing comparison neighborhoods, and let $L$ denote the Lebesgue number of the open cover $\{U\cap \gamma([0,1]) \colon U\in \mathcal{U}\}$ of $\gamma([0,1])$. Choose $\delta>0$ such that $d(\gamma(s),\gamma(t))<L$ for any $s,t\in [0,1]$ with $|t-s|<\delta$. Then picking points $0=t_0<\dots < t_N=1$ in such a way that $|t_{i+1}-t_i|<\delta/3$ for each $i$, we can guarantee that each segment $\gamma([t_i,t_{i+3}])$ lies entirely in one element of $\mathcal{U}$ ($i=0,\dots,N-3$). 

\begin{figure}[ht!] % [h!] is an optional placement specifier (here: "here, if possible!")
    \centering
    \includegraphics[width=0.8\linewidth,
                 height=0.5\textheight,
                 keepaspectratio]{conjv3.pdf}
    \caption{Geometric setup in the proof of Theorem \ref{Theorem: Timelike Rauch}}
    \label{fig:rauch-illustration} 
\end{figure}

Since the endpoints $p,q$ of $\gamma$ are symmetric conjugate, we may find sequences $(\gamma_n),(\sigma_n) \subseteq \TGeo(X)$ with $\gamma_n \neq \sigma_n$ for all $n$ but with the same endpoints $p_n,q_n$ and $\gamma_n,\sigma_n \to \gamma$ in the $d_{\Gamma}$-topology. It follows that there exist representatives (denoted in the same way) converging to $\gamma$ uniformly and in $\tau$-length. For $n$ large, we therefore also have that, for each $i$, the segments $\gamma_n([t_i,t_{i+3}])$ and $\sigma_n([t_i,t_{i+3}])$ lie in the same element of $\mathcal{U}$ as $\gamma([t_i,t_{i+3}])$.

In the remainder of the proof we additionally fix $n$ sufficiently large to secure that $\max(L_\tau(\gamma_n),L_\tau(\sigma_n))<D_\kappa$ and that $p_n\ll \gamma_n(t_1) \ll \sigma_n(t_2) \ll \gamma_n(t_3) \ll \dots \ll q_n$, as well as 
$p_n\ll \sigma_n(t_1) \ll \gamma_n(t_2) \ll \sigma_n(t_3) \ll \dots \ll q_n$ (cf.\ Figure \ref{fig:rauch-illustration}).
    

    We now consider the timelike triangle $\Delta(p_n,\gamma_n(t_1),\sigma_n(t_3))$ (with the obvious sides; for the side from $\gamma_n(t_1)$ to $\sigma_n(t_3)$ we choose a timelike maximizer entirely in the corresponding comparison neighborhood). This triangle can be subdivided into the smaller timelike triangles $\Delta(p_n,\gamma_n(t_1),\sigma_n(t_2))$ and $\Delta(\gamma_n(t_1),\sigma_n(t_2),\sigma_n(t_3))$. 
    
    
    This puts us into the situation of~\cite[Prop.\ 2.42]{beran2023splitting}, from which we conclude that the stacking of the corresponding comparison triangles of the sub-triangles in $\lm{\kappa}$ is concave. Analogously, we subdivide the next timelike triangle 
    $\Delta(\gamma_n(t_1),\sigma_n(t_3),\gamma_n(t_4))$ into the timelike sub-triangles $\Delta(\gamma_n(t_1),\gamma_n(t_2),\sigma_n(t_3))$ and $\Delta(\gamma_n(t_2),\sigma_n(t_3),\gamma_n(t_4))$ and iterate this procedure until we reach the endpoint $q_n$.
    In this way, we obtain piecewise timelike geodesics $\bar{\gamma}_n,\bar{\sigma}_n$ in $\lm{\kappa}$ with common endpoints whose Lorentzian arclength is the same as that of $\gamma_n,\sigma_n$, respectively. In particular, it is $< D_{\kappa}$. 

    We may assume that $\bar{\gamma}_n,\bar{\sigma}_n$ are parametrized to unit speed.
    Consider now a breakpoint $p$ of $\bar{\gamma}_n$, and for the moment call $\bar{\alpha}_n$ and $\bar{\beta}_n$ the geodesic sub-segments of $\bar{\gamma}_n$ meeting at $p$. Then replace $\bar{\beta}_n$ by its image under an isometry of 
    $\lm{\kappa}$ whose derivative maps the initial velocity vector of $\bar{\beta}_n$ to the terminal velocity vector of 
    $\bar{\alpha}_n$, thereby straightening the concatenation of $\bar{\alpha}_n$ and $\bar{\beta}_n$ into an unbroken geodesic
    segment. We repeat this procedure at each breakpoint of $\bar{\gamma}_n$, and analogously for $\bar{\sigma}_n$.
    Concavity of the stacking implies that these unbroken geodesics will remain within the region bounded by the original stacking and therefore also have to meet. Moreover, since we applied an isometry in each step, the resulting unbroken geodesic segments retain the Lorentzian length of $\bar{\gamma}_n$ and $\bar{\sigma}_n$, respectively. Consequently, we have found two different geodesic segments in $\lm{\kappa}$ that meet at $\bar{\tau}$-length strictly less than $D_{\kappa}$, a contradiction.
\end{proof}
\end{Theorem}


\subsection{Nonexistence of conjugate points under nonpositive timelike curvature}
\label{subsection: Conjugate points in spaces with timelike curvature bounds}

\begin{Theorem}
\label{Theorem: cartan hadamard adjacent}
Let $(X,d,\ll,\leq,\tau)$ be a Lorentzian pre-length space that is globally hyperbolic, strongly causal, locally causally closed, timelike path-connected,
and regular. Suppose that $X$ has curvature bounded above by $0$ in the sense of timelike triangle comparison and let $\gamma \in \TGeo(X)$. Then there exists a timelike family about $\gamma$ that is continuous at $\gamma$. This family is unique among timelike families about $\gamma$ that are continuous at $\gamma$. In particular, the endpoints of $\gamma$ are reachable (and are thus not ultimate conjugate).
\end{Theorem}
\begin{proof} By~\cite[Rem.\ 2.2]{Eroes-Gieger25}, $X$ is locally concave. Therefore, \cite[Thm.\ 3.1]{Eroes-Gieger25} shows, in particular, that there exists an open neighborhood $U$ of $\gamma([0,1])$ and open neighborhoods $U_0$ of $\gamma(0)$ and $U_1$ of $\gamma(1)$ with $U_0 \ll U_1$ such that the following properties hold:
\begin{itemize}
\item[(i)] For any $(p,q)\in U_0\times U_1$ there exists a unique timelike geodesic $\gamma_{pq} \colon [0,1]\to U$ with $\gamma_{pq}(0)=p$ and $\gamma_{pq}(1)=q$.
\item[(ii)] If $p_k\in U_0$ with $p_k\to \gamma(0)$ and $q_k \in U_1$ with $q_k \to q$, then $\gamma_{p_kq_k}$ 
converges uniformly to $\gamma$. In particular, also $d_F(\gamma_{p_kq_k},\gamma) \to 0$.
\end{itemize}
Cover $\gamma([0,1])$ by comparison neighborhoods $V_1,\dots,V_N$ such that for a suitable partition $0=t_0<\dots<t_N=1$, 
$\gamma([t_i,t_{i+1}])\subseteq V_{i+1}$ for all $i=0,\dots,N-1$. By (ii), there exists some $k_0$ such that also $\gamma_{p_kq_k}([t_i,t_{i+1}])\subseteq V_{i+1}$ for all $k\ge k_0$ and all $i$. The proof of~\cite[Prop.\ 2.3]{Eroes-Gieger25} shows that 
then each $\gamma_{p_kq_k}|_{[t_i,t_{i+1}]}$ as well as $\gamma|_{[t_i,t_{i+1}]}$ are timelike maximizers. Since $\tau|_{V_i\times V_i}$ is
continuous, we obtain
\begin{align*}
L_\tau(\gamma_{p_kq_k}|_{[t_i,t_{i+1}]}) = \tau(\gamma_{p_kq_k}(t_i),\gamma_{p_kq_k}(t_{i+1})) \to \tau(\gamma(t_i),\gamma(t_{i+1}))
= L_\tau (\gamma|_{[t_i,t_{i+1}]}).
\end{align*}
Consequently, 
\[
L_\tau(\gamma_{p_kq_k}) = \sum_{i=0}^{N-1} L_\tau(\gamma_{p_kq_k}|_{[t_i,t_{i+1}]}))
\to \sum_{i=0}^{N-1} L_\tau(\gamma|_{[t_i,t_{i+1}]})) = L_\tau(\gamma),
\]
implying that in fact $\gamma_{p_kq_k} \to \gamma$ with respect to $d_\Gamma$. 

It follows that $F \colon U_0\times U_1 \to \TGeo(X)$, $(p,q) \mapsto \gamma_{pq}$ is a timelike family about $\gamma$ that is
continuous at $\gamma$. If $\tilde F \colon W_0\times W_1 \to \TGeo(X)$ is another timelike family about $\gamma$ that is continuous at $\gamma$, then $d_\Gamma$-continuity of $\tilde F$ at $\gamma$ implies that there exist open neighborhoods $\tilde W_i
\subseteq W_i\cap U_i$ ($i=0,1$) such that $\tilde F(\tilde W_0\times \tilde W_1) \subseteq U$. Then (i) implies that 
$F=\tilde F$ on $\tilde W_0\times \tilde W_1$.

\end{proof}

\section{Outlook}
\label{section: Outlook}

In this work, we have introduced and studied timelike conjugate points in Lorentzian (pre-)length spaces. We have compared various notions with one another and with their Jacobi field counterpart in the smooth setting. Moreover, we have given applications (Theorems~\ref{Theorem: Timelike Rauch} and~\ref{Theorem: cartan hadamard adjacent}) illustrating their use.

There are several avenues for future work on conjugate points in the synthetic Lorentzian setting. First, a study of conjugate points along null geodesics would be of great interest. Here, care should be taken not to disrupt the compatibility with the smooth case. 

Moreover, a link between conjugate points and singularity theorems should be established. In this context, there are two tasks: Obtaining conjugate points from curvature bounds, and then using conjugate points to obtain singularities. A version of the latter was investigated by Minguzzi in the context of closed cone structures~\cite{minguzzi2019causality}.

\begin{appendices}

\section{The Fr\'echet distance between curves}
\label{Section: appendix frechet}

In this Appendix, we introduce the Fr\'echet distance between curves and prove some basic properties that were used in the main part of the paper. This notion indeed goes back to Fr\'echet and has found many applications (cf., e.g., \cite{Alt-Godau}), but we are not aware of a reference that rigorously establishes these fundamentals.

Let $(X,d)$ be a metric space. By a path in $X$ we mean an element $\gamma$ of the space of continuous maps $C([0,1],X)$ from $[0,1]$ into $X$ (we use the interval $[0,1]$ for convenience). Two paths $\gamma, \tilde{\gamma}$ are called equivalent if there exists a parametrization $\varphi \colon [0,1] \to [0,1]$ which is strictly increasing and surjective (hence continuous) such that $\gamma=\tilde{\gamma}\circ \varphi$. A curve is an equivalence class of paths under this equivalence relation, and we will
write $[\gamma]$ for the curve represented by the path $\gamma$. A path is called~\emph{nowhere constant\/} if it is not constant on any non-trivial sub-interval of $[0,1]$. This notion clearly descends to the space of curves as well. Given two curves $[\gamma],[\sigma]$, we define the~\emph{Fr\'echet distance\/} between them as follows:
\begin{equation}\label{eq:Frechet-distance}
    d_F([\gamma],[\sigma]) \coloneqq \inf_{\varphi,\psi} \max_{t \in [0,1]} d(\gamma(\varphi(t)),\sigma(\psi(t))),
\end{equation}
where $\varphi, \psi$ run over the set of all parametrizations as defined above. Obviously, we can equivalently write
\begin{equation*}
    d_F([\gamma],[\sigma]) \coloneqq \inf_{\varphi} \max_{t\in [0,1]} d(\gamma(\varphi(t)),\sigma(t)). 
\end{equation*}

We begin by establishing some properties of (spaces of) pairs of non-decreasing continuous maps:

\begin{Lemma}\label{Lem:parametrization} Let $\alpha \colon [0, 1] \to [0, 1]^2$ be a continuous, coordinate-wise non-decreasing map given by $\alpha(t) = (\varphi(t), \psi(t))$ with $\alpha(0) = (0, 0)$ and $\alpha(1) = (1, 1)$. Then there exists a coordinate-wise non-decreasing and $1$-Lipschitz continuous curve $\tilde{\alpha} = (\tilde{\varphi},\tilde{\psi}) \colon [0,2] \to [0,1]^2$ that has the same image as $\alpha$.
\end{Lemma}
\begin{proof} The function $s \colon [0,1]\to [0,2]$, $s(t) \coloneqq \varphi(t)+\psi(t)$ is continuous, non-decreasing, and surjective by the
intermediate value theorem. Let $\theta \colon [0,2] \to [0,1]$ be the generalized inverse of $s$,
\[
\theta(u) \coloneqq \inf \{ t\in [0,1] \mid s(t) \ge u \}.
\]
Since $s$ is surjective, $\theta$ is well-defined, and using continuity of $s$ one readily checks that $s(\theta(u)) = u$ for each $u\in [0,2]$. Now set, for $u\in [0,2]$:
\begin{align*}
    \tilde \alpha(u) \coloneqq \alpha(\theta(u)) = (\varphi(\theta(u)),\psi(\theta(u))) \eqqcolon (\tilde\varphi(u),\tilde\psi(u)).
\end{align*}
We first show that $\tilde\alpha([0,2]) = \alpha([0,1])$: Obviously, $\tilde\alpha([0,2]) \subseteq \alpha([0,1])$. Conversely, let
$p \coloneqq \alpha(t_0)$, $t_0\in [0,1]$ and $u_0 \coloneqq s(t_0)$. Setting $t_* \coloneqq \theta(s(t_0))$, we then have $\tilde\alpha(u_0) = \alpha(t_*)$. The definition of $\theta$, together with continuity and monotonicity of $s$ imply that $s(t_0)=s(t_*)$,
which we can also write as
\[
(\varphi(t_0)-\varphi(t_*)) + (\psi(t_0)-\psi(t_*)) = 0.
\]
Here, since $t_*\le t_0$ and since both $\varphi$ and $\psi$ are non-decreasing, both summands have to vanish, implying that $\alpha(t_0) = \alpha(t_*) =\tilde \alpha(u_0)$, establishing the claim.

Finally, to show that $\tilde \alpha$ is $1$-Lipschitz, let $0\le u<v\le 2$. Then using monotonicity of $\theta$, $\varphi$ and $\psi$,
hence of $\tilde \varphi$ and $\tilde\psi$,
\begin{align*}
\| \tilde\alpha(v) - \tilde \alpha(u)\|_1 &= |\tilde\varphi(v) - \tilde\varphi(u)| + |\tilde\psi(v)-\tilde\psi(u)|
= (\tilde\varphi(v) + \tilde\psi(v)) - (\tilde\varphi(u) + \tilde\psi(u))\\ 
&= s(\theta(v)) - s(\theta(u)) = v-u.
\end{align*}
Consequently, for any $u,v\in [0,2]$ we have
\[
\| \tilde\alpha(v) - \tilde \alpha(u)\|_2 \le \| \tilde\alpha(v) - \tilde \alpha(u)\|_1 \le |v-u|.
\]
\end{proof}
\begin{Lemma}\label{Lem:parametrizations-Hausdorff-compact}
Let $N$ be the set of continuous, coordinate-wise non-decreasing maps $\alpha=(\varphi,\psi) \colon [0,1]\to [0,1]^2$ with $\alpha(0) = (0,0)$
and $\alpha(1)=(1,1)$. Then $\mathcal{N} \coloneqq \{\alpha([0,1]) \mid \alpha\in N\}$ is compact with respect to the Hausdorff distance $d_H$.
\end{Lemma}
\begin{proof} Denote by $L$ the set of all coordinate-wise non-decreasing $1$-Lipschitz maps $\eta \colon [0,2]\to [0,1]^2$ with $\eta(0)=(0,0)$ and $\eta(1)=(1,1)$. By the Theorem of Arzel\`a--Ascoli, $L$ is relatively compact in $C([0,2],[0,1]^2)$. 
Since it is clearly also closed with respect to uniform convergence, it is indeed compact in $C([0,2],[0,1]^2)$.

Denote by $\mathcal{K}([0,1]^2)$ the space of compact subsets of $[0,1]^2$, and equip it with the Hausdorff distance $d_H$. Consider the map 
\begin{align*}
    \iota \colon C([0,2],[0,1]^2) &\to \mathcal{K}([0,1]^2)\\
    \eta &\mapsto \eta([0,2]).
\end{align*}
Since $d_H(\iota(\eta), \iota(\zeta)) \le \sup_{u\in [0,2]} \|\eta(u) - \zeta(u)\|_2 $, it follows that $\iota$ is continuous (indeed, $1$-Lipschitz). Thus $\mathcal{L} = \iota(L)$ is compact with respect to $d_H$. Finally, by Lemma~\ref{Lem:parametrization}, $\mathcal{N} = \mathcal{L}$, concluding the proof.
\end{proof}




Based on these preparations, we can now show:
\begin{Proposition}\label{Prop:Frechet-distance} The Fr\'echet distance $d_F$ defines a metric on the space of nowhere constant curves in $(X,d)$.  
\end{Proposition}
\begin{proof} Since parametrizations form a group, the Fr\'echet distance is well defined. Symmetry of $d_F$ is clear and the triangle inequality follows readily from the definition. Also $d_F([\gamma],[\gamma])=0$ is obvious. 

The only non-trivial point therefore is to show that $d_F([\gamma],[\sigma])=0$ implies $[\gamma]=[\sigma]$, i.e., that $\gamma$ is a reparametrization of $\sigma$. From the definition of $d_F$ we obtain sequences of parametrizations $(\varphi_n,\psi_n)$ such that, for each $n\in\N$:
\begin{equation}\label{eq:phin-psin}
    \max_{t\in [0,1]} d(\gamma(\varphi_n(t)),\sigma(\psi_n(t))) < \frac{1}{n}.
\end{equation}

Set $\alpha_n \coloneqq (\varphi_n,\psi_n)$, and $A_n \coloneqq \alpha_n([0,1])$. Then by Lemma~\ref{Lem:parametrizations-Hausdorff-compact}, $(A_n)_{n\in \N}$
is a sequence in the $d_H$-compact set $\mathcal{N}$, hence possesses a $d_H$-convergent subsequence, which we may without loss of generality assume to be $(A_n)_{n\in \N}$ itself. Let $A\in \mathcal{N}$ be its limit. Then there is some $\alpha = (\varphi,\psi)\in N$ with $A=\alpha([0,1])$. Given any $t\in [0,1]$, there exists a sequence $(t_n)_{n\in \N}$ in $[0,1]$ such that 
\[
A_n\ni (\varphi_n(t_n),\psi_n(t_n)) \to (\varphi(t),\psi(t)) \qquad (n\to \infty).
\]
By \eqref{eq:phin-psin} we have $d(\gamma(\varphi_n(t_n)),\sigma(\psi_n(t_n))) < \frac{1}{n}$, so letting $n\to \infty$ we obtain
$d(\gamma(\varphi(t)),\sigma(\psi(t)))=0$, implying that $\gamma\circ \varphi = \sigma\circ \psi$. 

We are left with the task of constructing from $\varphi$ and $\psi$ new continuous surjective maps $\hat\varphi$, $\hat\psi \colon [0,1]\to [0,1]$ that are strictly monotonically increasing, hence qualify as parametrizations and satisfy $\gamma\circ \hat\varphi = \sigma\circ \hat\psi$. This we do
in several steps.

We first claim that $\varphi$ is constant on any non-trivial interval $[t_1,t_2]$ of $[0,1]$ if and only if $\psi$ is constant on $[t_1,t_2]$.
By symmetry, it suffices to show the `only if' part of the claim. Since $\varphi$ is constant on $[t_1,t_2]$, so is
$\gamma\circ \varphi = \sigma\circ \psi$.
Also, $\psi$ being continuous and non-decreasing, the image of $[t_1, t_2]$ under $\psi$ is the closed interval $J = [\psi(t_1), \psi(t_2)]$.
The path $\sigma$ is therefore constant on $J$. By our assumption on $\sigma$ this forces $J$ to be a single point, hence
$\psi(t_1) = \psi(t_2)$. Since $\psi$ is non-decreasing, $\psi$ is indeed constant on $[t_1, t_2]$.

Define $\tau \colon [0,1] \to [0,1]$ by
$$
\tau(t) = \frac{\varphi(t) + \psi(t)}{2}.
$$
Then $\tau$ is continuous, non-decreasing and surjective $[0,1]\to [0,1]$. 

Next we show that $\tau(t)$ is constant on an interval $[t_1, t_2]$ if and only if both $\vp(t)$ and $\ps(t)$ are constant on that interval. For the non-trivial direction of this claim, suppose that $\tau$ is constant on $[t_1, t_2]$. Then for any $t, t'$ with $t_1 \le t < t' \le t_2$, we have $\tau(t) = \tau(t')$, hence
\[
\frac{\vp(t) + \ps(t)}{2} = \frac{\vp(t') + \ps(t')}{2} \implies (\vp(t') - \vp(t)) + (\ps(t') - \ps(t)) = 0
\]
Since $\vp$ and $\ps$ are non-decreasing, we conclude that $\vp(t) = \vp(t')$ and $\ps(t) = \ps(t')$. Thus both $\vp$ and $\ps$ are constant on $[t_1, t_2]$.


For any $s \in [0,1]$, define $\hvp(s) \coloneqq \vp(t)$ and $\hps(s) \coloneqq \psi(t)$ for any $t \in \tau^{-1}(s)$. Since $\vp$, $\psi$ are constant on the fibers $\tau^{-1}(s)$ of $\tau$, we obtain well-defined functions 
$\hvp \colon [0,1] \to [0,1]$ and $\hps \colon [0,1] \to [0,1]$ such that
\[
\vp(t) = \hvp(\tau(t)) \quad \text{and} \quad \ps(t) = \hps(\tau(t)) \quad \text{for all } t \in [0,1].
\]
By construction, $\g\circ\vp = \sigma\circ \psi$, so also $\g(\hvp(\tau(t))) = \s(\hps(\tau(t)))$ for all $t\in [0,1]$.
Surjectivity of $\tau$ then gives $\g\circ \hvp = \sigma\circ \hps$. Also, since $\varphi$, $\psi$ are surjective, so are $\hvp$ and $\hps$.

Since $\tau$ is a continuous map from a compact space ($[0,1]$) to a Hausdorff space (again $[0,1]$), it is a quotient map. By the universal property of quotient maps, since $\vp$ and $\ps$ are continuous and constant on the fibers of $\tau$, the induced maps $\hvp$ and $\hps$ are also continuous.

It remains to show strict monotonicity, which we do for $\hat\varphi$. Let $0 \le s_1 < s_2 \le 1$. We must show $\hvp(s_1) < \hvp(s_2)$.
Since $\tau$ is non-decreasing and surjective from $[0,1]$ to $[0,1]$, we can find $t_1, t_2 \in [0,1]$ such that $\tau(t_1) = s_1$, $\tau(t_2) = s_2$, and $t_1 < t_2$. Then
\[
\hat\varphi(s_1) = \varphi(t_1) \le \varphi(t_2) = \hat\varphi(s_2).
\]
Let us indirectly assume that equality holds: $\hvp(s_1) = \hvp(s_2)$. Then $\vp(t_1) = \vp(t_2)$ and since $\vp$ is non-decreasing, it must be constant on the entire interval $[t_1, t_2]$. By what was shown above, then $\ps$ must also be constant on $[t_1, t_2]$.
But then also $\tau(t)$ must be constant on $[t_1, t_2]$. In particular, $\tau(t_1) = \tau(t_2)$, which implies $s_1 = s_2$, a contradiction.
This proves that $\hvp$ is strictly monotonically increasing, as claimed.
\end{proof}
\end{appendices}

\section*{Acknowledgments}

The authors would like to thank Christina Sormani and Nathalie Tassotti for helpful discussions.

This research was funded in part by the Austrian Science Fund (FWF) [Grant DOI 10.55776/EFP6 and 10.55776/J4913]. For open access purposes, the authors have applied a CC BY public copyright license to any author accepted manuscript version arising from this submission.

%\bibliography{bibliography} 
%\bibliographystyle{acm}
%\addcontentsline{toc}{section}{References}
%\bibliographystyle{alpha}

\begin{thebibliography}{10}

\bibitem{Alt-Godau}
{\sc Alt, H., and Godau, M.}
\newblock Computing the {Fr{\'e}chet} distance between two polygonal curves.
\newblock {\em Int. J. Comput. Geom. Appl. 5}, 1-2 (1995), 75--91.

\bibitem{BKR24}
{\sc Beran, T., Kunzinger, M., and Rott, F.}
\newblock On curvature bounds in {Lorentzian} length spaces.
\newblock {\em J. Lond. Math. Soc., II. Ser. 110}, 2 (2024), 41.
\newblock Id/No e12971.

\bibitem{beran2023splitting}
{\sc Beran, T., Ohanyan, A., Rott, F., and Solis, D.~A.}
\newblock The splitting theorem for globally hyperbolic {Lorentzian} length
 spaces with non-negative timelike curvature.
\newblock {\em Letters in Mathematical Physics 113}, 2 (2023), 48.

\bibitem{beran2022angles}
{\sc Beran, T., and S{\"a}mann, C.}
\newblock Hyperbolic angles in {L}orentzian length spaces and timelike
 curvature bounds.
\newblock {\em J.\ London Math.\ Soc.\ 2\/} (2023), 1--58.

\bibitem{bridson2013metric}
{\sc Bridson, M.~R., and Haefliger, A.}
\newblock {\em Metric spaces of non-positive curvature}, vol.~319.
\newblock Springer Science \& Business Media, 2013.

\bibitem{EVP21}
{\sc Ebrahimi, N., Vatandoost, M., and Pourkhandani, R.}
\newblock On conformal {L}orentzian length spaces.
\newblock {\em Anal. Math. Phys. 13}, 6 (2023), Paper No. 93, 20.

\bibitem{Eroes-Gieger25}
{\sc Er{\"o}s, D., and Gieger, S.}
\newblock A synthetic {Lorentzian} {Cartan}-{Hadamard} theorem.
\newblock Transactions Amer.\ Math.\ Soc., to appear, {arXiv}:2506.22197 [math.{MG}].

\bibitem{flores2024low}
{\sc Flores, J.~L., Herrera, J., and Solis, D.~A.}
\newblock Low regularity approach to bartnik’s conjecture.
\newblock {\em Classical and Quantum Gravity\/} (2024).

\bibitem{grant2019inextendibility}
{\sc Grant, J.~D.~E., Kunzinger, M., and S{\"a}mann, C.}
\newblock Inextendibility of spacetimes and {Lorentzian} length spaces.
\newblock {\em Annals of global analysis and geometry 55}, 1 (2019), 133--147.

\bibitem{GMST:05}
{\sc Gravesen, J., Markvorsen, S., Sinclair, R., Tanaka, M.}
\newblock {The cut locus of a torus of revolution.}
\newblock {\em Asian J.\ Math.\ 9}, 1 (2005), 103--120.

\bibitem{hau2020causal}
{\sc Hau, L.~A., Pacheco, A.~J.~C., and Solis, D.~A.}
\newblock On the causal hierarchy of {Lorentzian} length spaces.
\newblock {\em Classical and Quantum Gravity 37}, 21 (2020), 215013.

\bibitem{KOSS}
{\sc Kunzinger, M., Ohanyan, A., Schinnerl, B., and Steinbauer, R.}
\newblock The {H}awking-{P}enrose singularity theorem for {$C^1$}-{L}orentzian
 metrics.
\newblock {\em Comm. Math. Phys. 391}, 3 (2022), 1143--1179.

\bibitem{kunzinger2018lorentzian}
{\sc Kunzinger, M., and S{\"a}mann, C.}
\newblock Lorentzian length spaces.
\newblock {\em Annals of global analysis and geometry 54}, 3 (2018), 399--447.

\bibitem{kunzinger2014regularisation}
{\sc Kunzinger, M., Steinbauer, R., Stojkovic, M., and Vickers, J.~A.}
\newblock A regularisation approach to causality theory for
 {$C^{1,1}$}-{Lorentzian} metrics.
\newblock {\em General Relativity and Gravitation 46}, 8 (2014), 1738.

\bibitem{McCann_Saemann22}
{\sc McCann, R.~J., and S\"amann, C.}
\newblock A {L}orentzian analog for {H}ausdorff dimension and measure.
\newblock {\em Pure Appl. Anal. 4}, 2 (2022), 367--400.

\bibitem{minguzzi2015convex}
{\sc Minguzzi, E.}
\newblock Convex neighborhoods for {Lipschitz} connections and sprays.
\newblock {\em Monatshefte f{\"u}r Mathematik 177}, 4 (2015), 569--625.

\bibitem{minguzzi2019causality}
{\sc Minguzzi, E.}
\newblock Causality theory for closed cone structures with applications.
\newblock {\em Reviews in Mathematical Physics 31}, 05 (2019), 1930001.

\bibitem{ON83}
{\sc O'Neill, B.}
\newblock {\em Semi-{R}iemannian geometry}, vol.~103 of {\em Pure and Applied
 Mathematics}.
\newblock Academic Press, Inc. [Harcourt Brace Jovanovich, Publishers], New
 York, 1983.
\newblock With applications to relativity.

\bibitem{rinow1961innere}
{\sc Rinow, W.}
\newblock {\em Die innere Geometrie der metrischen R{\"a}ume}, vol.~105.
\newblock Springer, 1961.

\bibitem{samann2018geodesics}
{\sc S{\"a}mann, C., and Steinbauer, R.}
\newblock On geodesics in low regularity.
\newblock In {\em Journal of Physics: Conference Series\/} (2018), vol.~968,
 IOP Publishing, p.~012010.

\bibitem{shankar2009conjugate}
{\sc Shankar, K., and Sormani, C.}
\newblock Conjugate points in length spaces.
\newblock {\em Advances in Mathematics 220}, 3 (2009), 791--830.

\bibitem{szeghy2008conjugate}
{\sc Szeghy, D.}
\newblock On the conjugate locus of pseudo-{R}iemannian manifolds.
\newblock {\em Indagationes Mathematicae 19}, 3 (2008), 465--480.

\bibitem{Treude:2011}
{\sc Treude, J.-H.}
\newblock Ricci curvature comparison in {Riemannian} and {Lorentzian} geometry.
\newblock Master's thesis, Universität Freiburg, 2011.

\bibitem{warner1965conjugate}
{\sc Warner, F.~W.}
\newblock The conjugate locus of a {R}iemannian manifold.
\newblock {\em American Journal of Mathematics 87}, 3 (1965), 575--604.

\end{thebibliography}

\end{document}
